\documentclass[11pt,letterpaper,oneside,openany]{book}
\usepackage[T1]{fontenc}
\usepackage[utf8]{inputenc}
\usepackage{amsmath,amssymb,mathtools,amsthm}
% Canonical mathematical notation for active FabiusFunction LaTeX documents.
%
% This file is intentionally a plain TeX input rather than a LaTeX package.
% Every standalone document loads it by an explicit path relative to that
% document's directory; included fragments inherit the definitions from their
% root document.  Keep this file independent of document layout and metadata.
%
% Contract:
%   * one command has one mathematical meaning throughout the active corpus;
%   * names encode semantic roles rather than a document-local choice of letter;
%   * visually similar but differently normalized objects have different print;
%   * every command below is defined normatively in the notation catalogue.
%
% The active corpus excludes docs/archive/.  Historical sources may retain
% their original notation only inside an explicitly labelled quotation.

\makeatletter
\@ifundefined{mathbb}{%
  \PackageError{fabius-notation}{The amssymb package must be loaded first}%
    {Load amsmath and amssymb before inputting fabius-notation.tex.}%
}{}
\@ifundefined{operatorname}{%
  \PackageError{fabius-notation}{The amsmath package must be loaded first}%
    {Load amsmath before inputting fabius-notation.tex.}%
}{}
\makeatother

% Version stamp used by the catalogue and by static audits.
\newcommand{\FabiusNotationVersion}{2026-08-31}

% Number systems.  NaturalNumbers is explicitly the nonnegative convention.
\newcommand{\NaturalNumbers}{\mathbb{N}_{0}}
\newcommand{\PositiveIntegers}{\mathbb{N}_{>0}}
\newcommand{\IntegerNumbers}{\mathbb{Z}}
\newcommand{\RationalNumbers}{\mathbb{Q}}
\newcommand{\RealNumbers}{\mathbb{R}}
\newcommand{\ComplexNumbers}{\mathbb{C}}
\newcommand{\ExtendedRealNumbers}{\overline{\mathbb{R}}}

% The three principal Fabius--Rvachev functions.  Subscripts are deliberate:
% a bare F and a calligraphic F have several unrelated uses in the corpus.
\newcommand{\FabiusBounded}{F_{\mathrm{bd}}}
\newcommand{\FabiusGlobal}{F_{\mathrm{gl}}}
\newcommand{\FabiusQuantile}{Q_{F}}
\newcommand{\FabiusClampedQuantile}{Q_{F,\mathrm{cl}}}
\newcommand{\RvachevUp}{\operatorname{up}}

% Constants, elementary operators, and delimiters.
\newcommand{\EulerE}{\mathrm{e}}
\newcommand{\ImaginaryUnit}{\mathrm{i}}
\newcommand{\Differential}{\mathop{}\!\mathrm{d}}
\newcommand{\AbsoluteValue}[1]{\left\lvert #1\right\rvert}
\newcommand{\Norm}[1]{\left\lVert #1\right\rVert}
\newcommand{\Floor}[1]{\left\lfloor #1\right\rfloor}
\newcommand{\Ceiling}[1]{\left\lceil #1\right\rceil}
\newcommand{\FractionalPart}[1]{\left\{#1\right\}}
\newcommand{\PositivePart}[1]{\left(#1\right)_{+}}
\newcommand{\IndicatorOf}[1]{\mathbf{1}_{#1}}
\newcommand{\SetOf}[1]{\left\{#1\right\}}
\newcommand{\InnerProduct}[1]{\left\langle #1\right\rangle}
\newcommand{\CardinalityOf}[1]{\left\lvert #1\right\rvert}
\newcommand{\DefinitionEquals}{\mathrel{\coloneqq}}
\newcommand{\Sign}{\operatorname{sgn}}
\newcommand{\IdentityOperator}{\operatorname{id}}
\newcommand{\SpanOperator}{\operatorname{span}}
\newcommand{\TraceOperator}{\operatorname{tr}}
\newcommand{\RankOperator}{\operatorname{rank}}
\newcommand{\DiagonalOperator}{\operatorname{diag}}
\newcommand{\ResidueOperator}{\operatorname*{Res}}
\newcommand{\LeastCommonMultiple}{\operatorname{lcm}}
\newcommand{\OrderOperator}{\operatorname{ord}}
\newcommand{\PrincipalLogarithm}{\operatorname{Log}}
\newcommand{\FallingFactorial}[2]{\left(#1\right)^{\underline{#2}}}
\newcommand{\RisingFactorial}[2]{\left(#1\right)^{\overline{#2}}}

% Support, topology, convolution, and metric notation.
\newcommand{\SupportOperator}{\operatorname{supp}}
\newcommand{\SupportOf}[1]{\SupportOperator\!\left(#1\right)}
\newcommand{\TopologicalSupportOf}[1]{\operatorname{tsupp}\!\left(#1\right)}
\newcommand{\Convolution}{\mathbin{*}}
\newcommand{\MetricDistance}{\operatorname{dist}}
\newcommand{\TotalVariationDistance}{d_{\mathrm{TV}}}

% Probability.  Equality and convergence in law are not metric distance.
\newcommand{\Probability}{\mathbb{P}}
\newcommand{\Expectation}{\mathbb{E}}
\newcommand{\Variance}{\operatorname{Var}}
\newcommand{\Covariance}{\operatorname{Cov}}
\newcommand{\LawOperator}{\operatorname{Law}}
\newcommand{\LawOf}[1]{\LawOperator\!\left(#1\right)}
\newcommand{\EqualInLaw}{\mathrel{\overset{\mathrm{law}}{=}}}
\newcommand{\ConvergesInLaw}{\mathrel{\xrightarrow{\mathrm{law}}}}

% Fourier and integral transforms.  The printed labels expose normalization.
% FourierTwoPi uses exp(-2*pi*i*x*xi); FourierAngular uses exp(-i*x*omega).
\newcommand{\FourierTwoPi}[1]{\widehat{#1}_{\,2\pi}}
\newcommand{\FourierAngular}[1]{\widehat{#1}_{\,\mathrm{ang}}}
\newcommand{\LaplaceTransformSymbol}{\mathcal{L}}
\newcommand{\MellinTransformSymbol}{\mathcal{M}}
\newcommand{\LaplaceTransformOf}[1]{\LaplaceTransformSymbol\!\left[#1\right]}
\newcommand{\MellinTransformOf}[1]{\MellinTransformSymbol\!\left[#1\right]}
\newcommand{\MomentGeneratingFunction}[1]{M_{#1}}
\newcommand{\CharacteristicFunction}[1]{\varphi_{#1}}

% The two standard sinc functions must never share one symbol.
\newcommand{\SincRad}{\operatorname{sinc}_{\mathrm{rad}}}
\newcommand{\SincPi}{\operatorname{sinc}_{\pi}}

% Binary arithmetic and Thue--Morse notation.
\newcommand{\BinaryDigitSum}{\operatorname{wt}_{2}}
\newcommand{\ThueMorseSignSymbol}{\varepsilon}
\newcommand{\ThueMorseSign}[1]{\varepsilon_{#1}}
\newcommand{\PadicValuation}[1]{\nu_{#1}}
\newcommand{\TwoAdicValuation}{\nu_{2}}

% q-series notation.  Every base is an explicit argument.
\newcommand{\QPochhammer}[3]{\left(#1;#2\right)_{#3}}
\newcommand{\GaussianBinomial}[3]{%
  \genfrac{[}{]}{0pt}{}{#1}{#2}_{#3}%
}
\newcommand{\QInteger}[2]{\left[#1\right]_{#2}}

% Named special functions.
\newcommand{\Polylogarithm}{\operatorname{Li}}
\newcommand{\LambertW}{W}
\newcommand{\GammaFunction}{\Gamma}
\newcommand{\RiemannZeta}{\zeta}

% Asymptotics.  BigO and LittleO are operator tokens for existing formula
% grammar; prefer the At forms so that the limiting regime is printed.
\newcommand{\BigO}{\mathcal{O}}
\newcommand{\LittleO}{o}
\newcommand{\BigOAt}[2]{\mathcal{O}_{#2}\!\left(#1\right)}
\newcommand{\LittleOAt}[2]{o_{#2}\!\left(#1\right)}

\IfFileExists{libertinus.sty}{\usepackage{libertinus}}{\usepackage{lmodern}}
\usepackage[final]{microtype}
\usepackage[letterpaper,margin=0.90in,headheight=15pt]{geometry}
\usepackage{bm}
\usepackage{booktabs,longtable,tabularx,array,multirow}
\usepackage{enumitem}
\usepackage{xcolor}
\usepackage{fancyhdr}
\usepackage{graphicx}
\usepackage{float}
\usepackage{caption}
\usepackage{mathrsfs}
\usepackage{listings}
\usepackage{xurl}
\usepackage{hyperref}
\usepackage[nameinlink,capitalise,noabbrev]{cleveref}
\usepackage{bookmark}

\definecolor{linkblue}{RGB}{18,74,130}
\definecolor{citegreen}{RGB}{24,104,74}
\definecolor{provedgreen}{RGB}{20,100,60}
\definecolor{conjectureorange}{RGB}{158,82,0}
\definecolor{numericalblue}{RGB}{28,78,135}
\definecolor{shadegray}{RGB}{246,247,249}
\hypersetup{
  colorlinks=true,
  linkcolor=linkblue,
  citecolor=citegreen,
  urlcolor=linkblue,
  pdftitle={An Inverse-Jet Atlas for q-Analogs},
  pdfsubject={Multivariate reversion and singular inverse expansions for q-Pochhammer symbols, Gaussian coefficients, q-gamma, q-beta, and related q-analogs},
  pdfkeywords={inverse q-Pochhammer, inverse q-binomial, Puiseux series, roots of unity, q-gamma, q-beta, Lagrange inversion, Fabius function}
}
\allowdisplaybreaks[3]
\setcounter{tocdepth}{2}
\setcounter{secnumdepth}{3}
\setlist{itemsep=2pt,topsep=5pt,leftmargin=2.2em}
\setlength{\emergencystretch}{3em}

\pagestyle{fancy}
\fancyhf{}
\fancyhead[L]{\nouppercase{\leftmark}}
\fancyhead[R]{\thepage}
\renewcommand{\headrulewidth}{0.4pt}

\newtheorem{theorem}{Theorem}[chapter]
\newtheorem{lemma}{Lemma}[chapter]
\newtheorem{proposition}{Proposition}[chapter]
\newtheorem{corollary}{Corollary}[chapter]
\newtheorem{principle}{Principle}[chapter]
\newtheorem{identity}{Identity}[chapter]
\theoremstyle{definition}
\newtheorem{definition}{Definition}[chapter]
\newtheorem{example}{Example}[chapter]
\newtheorem{algorithm}{Algorithm}[chapter]
\newtheorem{researchproblem}{Research problem}[chapter]
\newtheorem{conjecture}{Conjecture}[chapter]
\theoremstyle{remark}
\newtheorem{remark}{Remark}[chapter]
\newtheorem{computationalresult}{Computational result}[chapter]
\newtheorem{warning}{Warning}[chapter]

% Some TeX distributions accumulate theorem-counter reset chains while
% independent theorem families are declared in succession.  Each counter is
% already reset directly by the chapter counter; clearing their subordinate
% reset lists prevents accidental cycles without changing the numbering.
\makeatletter
\let\cl@theorem\@empty
\let\cl@lemma\@empty
\let\cl@proposition\@empty
\let\cl@corollary\@empty
\let\cl@principle\@empty
\let\cl@identity\@empty
\let\cl@definition\@empty
\let\cl@example\@empty
\let\cl@algorithm\@empty
\let\cl@researchproblem\@empty
\let\cl@conjecture\@empty
\let\cl@remark\@empty
\let\cl@computationalresult\@empty
\let\cl@warning\@empty
\makeatother

\newcommand{\Disc}{\operatorname{Disc}}
\newcommand{\ord}{\operatorname{ord}}
\newcommand{\qmulti}[2]{\genfrac{[}{]}{0pt}{}{#1}{#2}_{q}}
\newcommand{\Aq}{\mathcal A_q}
\newcommand{\PhiC}{\Phi}
\newcommand{\statusproved}{\textcolor{provedgreen}{\textsc{proved}}}
\newcommand{\statusnumerical}{\textcolor{numericalblue}{\textsc{symbolic/numerical}}}
\newcommand{\statusconjectural}{\textcolor{conjectureorange}{\textsc{conjectural}}}
\newcommand{\ind}{\mathbf{1}}

\lstdefinestyle{pythonstyle}{
  language=Python,
  basicstyle=\ttfamily\footnotesize,
  backgroundcolor=\color{shadegray},
  frame=single,
  breaklines=true,
  columns=fullflexible,
  keepspaces=true,
  showstringspaces=false
}

\title{\Huge An Inverse-Jet Atlas for \(q\)-Analogs\\[0.45em]
\Large Multivariate Reversion, Singular Branches, and Asymptotics at\\
\(q=0,1,-1\), Roots of Unity, and \(q=\pm\infty\)}
\author{A second research companion to the ProveIt \(q\)-series drafts}
\date{August 30, 2026}

\begin{document}
\frontmatter
\hypersetup{pageanchor=false}
\maketitle
\clearpage
\hypersetup{pageanchor=true}

\chapter*{Abstract}
\addcontentsline{toc}{chapter}{Abstract}
The ProveIt draft corpus already contains a substantial forward theory of
\(q\)-Pochhammer symbols, Gaussian coefficients, and their asymptotic
expansions, together with a first report on inverse \(q\)-analogs.  The present
volume deliberately starts beyond that baseline.  Its organizing object is an
\emph{inverse jet}: a local inverse expansion equipped with its recovered
parameter, base point, ramification index, branch label, domain of validity,
and dependence on every remaining parameter.  Ordinary Taylor inverses,
Puiseux inverses, logarithmic inverses, and multiscale inverses are treated in
one framework.

Several results appear to be new within the audited repository corpus.  For the
finite product \(P_n(a,q)=\QPochhammer{a}{q}{n}\), the nontrivial parameter for which
\(q=-1\) becomes critical is exactly
\[
 a_* = \frac1{n-1} \quad(n\text{ even}),
 \qquad
 a_*=-\frac1n \quad(n\text{ odd}).
\]
For every \(n\ge4\) this produces a quadratic fold and hence two square-root
inverse branches.  The sole low-order exception is the exact cubic identity
\[
 \QPochhammer{-1/3}{q}{3}=\frac{32}{27}+\frac4{27}(q+1)^3,
\]
which gives three cube-root branches without a remainder.  At arbitrary roots
of unity we express the ramification order by cyclotomic multiplicity.  A
Gaussian polynomial is square-free because each multiplicity is \(0\) or
\(1\); a \(q\)-multinomial can have higher multiplicity, equal to an explicit
residue-carry count.  In particular,
\[
 [n]_q! = \lfloor n/2\rfloor!\,(q+1)^{\lfloor n/2\rfloor}
       (1+\BigO(q+1))
\]
near \(q=-1\), so its inverse has a precisely normalized higher Puiseux branch.

For the infinite product we derive a fixed-target double-scaling inverse that
is distinct from the usual fixed-argument \(q\to1\) expansion.  If
\((a;\EulerE^{-\varepsilon})_\infty=\EulerE^{-L}\), with \(L>0\) fixed, then
\[
\begin{aligned}
 a={}&L\varepsilon-\frac{L(L+2)}4\varepsilon^2
 +\frac{L(L^2+9L+12)}{72}\varepsilon^3 \\
 &-\frac{L(L^3+4L^2+12L+24)}{576}\varepsilon^4
 +\BigO(\varepsilon^5).
\end{aligned}
\]
An arbitrary-order recurrence is supplied.  For \(q\)-gamma we obtain an
all-order algebraic expansion in which every odd power beyond the linear term
vanishes.  Reversion yields expansions both for the base and for the argument.
For \(q\)-beta, the base inverse changes type on the explicit critical curve
\(xy=1\): it is ordinary off the curve, square-root on the curve except at the
identically constant point \((1,1)\).

For Gaussian coefficients we prove a general centered-palindromic theorem.  A
nonconstant polynomial with nonnegative palindromic coefficients becomes, after
centering, the moment-generating function of a symmetric finite distribution.
Its normalized inverse therefore has exactly two global real branches, and its
branch at \(q=1\) is governed by cumulants.  This covers Gaussian and
\(q\)-multinomial coefficients and many finite product analogs.  The report also
treats inverses of \(q\)-numbers, \(q\)-factorials, complex orders,
\(q\)-gamma, \(q\)-digamma, \(q\)-beta, Euler \(q\)-exponentials, basic
hypergeometric series, theta products, and the effective geometric base in
Fabius--Rvachev formulas.  Proved results, reproducible symbolic/numerical
checks, and conjectures are explicitly separated.

\chapter*{Corpus audit and novelty boundary}
\addcontentsline{toc}{chapter}{Corpus audit and novelty boundary}
The source snapshot is the \texttt{main} branch of Vladimir Reshetnikov's
\texttt{ProveIt} repository at commit
\texttt{1cea73234a363ddbc392816f6babb5a57920e984}, dated 30 August 2026.
The live subtree named in the research request contains the large synthesis
\emph{Exponents and \(q\)-Series Frontiers}, the Fabius parameter-response and
Rvachev frontier reports, the forward expansion report, the proof-oriented
\(q\)-Pochhammer/\(q\)-binomial monograph, and the earlier inverse report,
together with small generated TeX data tables.  The code-search index lagged
behind the live Contents tree, so the audit used the latter and the current
commit tree rather than the stale search count.

The following material is treated as established repository prior art and is
not repackaged as novelty here:
\begin{enumerate}
\item the finite/infinite product identities, Gaussian recurrences,
      \(q\)-gamma interpolation, and Bernoulli/polylogarithmic forward
      expansions developed in the monograph and expansion report;
\item the first inverse report's real branch classifications, principal inverse
      of \((a;q)_\infty\), positive-base uniqueness for Gaussian coefficients,
      secondary discriminants, and numerical branch tracking;
\item the Fabius--Rvachev use of the fixed half-base \(q=1/2\), geometric
      coefficient arrays, infinite sinc products, Bernoulli--Bell moment
      machinery, and parameter-response constructions.
\end{enumerate}
The new emphasis is the missing \emph{inverse atlas}: explicit inverse jets at
all singular bases, systematic inversion in every parameter, and transport
between regular and ramified regimes.  ``New'' means not located in the active
corpus under equivalent notation; it is not an unconditional claim of
bibliographic priority.

\begin{center}
\begin{tabularx}{0.96\textwidth}{@{}lX@{}}
\toprule
Label & Meaning \\
\midrule
\statusproved & A proof is given here, or a precisely cited standard theorem is
used with its hypotheses displayed. \\
\statusnumerical & An exact computer-algebra identity or high-precision check is
reproduced by the accompanying Python program. \\
\statusconjectural & A precise conjecture or research direction; no proof is
claimed. \\
\bottomrule
\end{tabularx}
\end{center}

\chapter*{How to read the atlas}
\addcontentsline{toc}{chapter}{How to read the atlas}
An entry for an inverse branch records six pieces of data:
\[
 (F,p_0,y_0;m;\mathcal B;\mathcal S).
\]
Here \(F\) is the forward function; \(p\) is the recovered parameter;
\((p_0,y_0)\) is the base point; \(m\) is the ramification order; \(\mathcal
B\) is the branch convention; and \(\mathcal S\) is the singular set in all
auxiliary parameters.  Thus \(m=1\) means an ordinary Taylor inverse,
\(m>1\) a Puiseux inverse, and \(m=\infty\) is mnemonic for an essential or
logarithmic inversion.  An inverse without these data is normally ambiguous.

The variable
\[
 q=\EulerE^{-\varepsilon}
\]
is used near \(q=1\).  It preserves reciprocal symmetry: \(q\mapsto q^{-1}\)
becomes \(\varepsilon\mapsto-\varepsilon\).  Near infinity we use
\(Q=q^{-1}\).  At a primitive \(d\)-th root of unity we write \(q=\zeta+h\).
The target is frequently logarithmically normalized:
\[
 w=\log\frac{y}{F(p_0)},
\]
which converts products into additive cumulant jets.

\tableofcontents

\mainmatter

\part{Universal inversion machinery}

\chapter{Parameter-wise inverses and their geometry}

\section{One equation does not recover every parameter}
Let \(F(p_1,\ldots,p_r)=y\) be a scalar \(q\)-analog.  Unless additional
relations are supplied, the level set is generically an \((r-1)\)-dimensional
hypersurface.  Consequently an ``inverse of \(F\)'' always means either a
partial inverse in one named parameter or a simultaneous inverse of a square
system.

\begin{definition}[Partial inverse branch]
Fix auxiliary parameters \(\lambda\) and write \(F(p;\lambda)\).  A partial
inverse branch is a function \(p=G(y;\lambda)\), together with a base point and
continuation rule, such that
\[
 F(G(y;\lambda);\lambda)=y.
\]
Its singular values are the images of points where \(F_p=0\), together with
values induced by poles, essential singularities, or singular auxiliary
parameters.
\end{definition}

\begin{theorem}[\statusproved\ Local partial inverse and first sensitivities]
Suppose \(F\) is holomorphic near \((p_0,\lambda_0)\), put
\(y_0=F(p_0,\lambda_0)\), and assume \(F_p(p_0,\lambda_0)\ne0\).  Then there is
a unique local holomorphic branch \(p=G(y;\lambda)\) with
\(G(y_0;\lambda_0)=p_0\).  On this branch,
\[
 G_y=\frac1{F_p},
 \qquad
 G_{\lambda_j}=-\frac{F_{\lambda_j}}{F_p}.
\]
\end{theorem}
\begin{proof}
The holomorphic implicit-function theorem gives existence and uniqueness.
Differentiate \(F(G(y;\lambda),\lambda)=y\) first in \(y\) and then in
\(\lambda_j\).
\end{proof}

These identities are also condition-number formulas.  A relative condition
number for recovering \(p\) from \(y\) is
\[
 \kappa_{p\leftarrow y}=
 \AbsoluteValue{\frac{y}{pF_p}},
\]
when \(pyF_p\ne0\).  It diverges at a fold and therefore predicts the
transition from Taylor to Puiseux inversion.

\section{The inverse Riemann surface}
For a fixed auxiliary parameter the analytic curve
\[
 \mathscr R_F=\{(p,y)\in\ComplexNumbers^2:F(p)=y\}
\]
projects to the target plane.  Away from critical values the projection is a
covering.  Loops in the punctured target plane permute the sheets; this
monodromy is part of the inverse, not an implementation detail.  For a finite
\(q\)-Pochhammer product inverted in its argument, the covering has \(n\)
sheets.  Inverting in \(q\) gives degree \(n(n-1)/2\) generically.  Infinite
products replace the algebraic covering by an infinite-sheeted analytic one.

\section{Branch labels used below}
Real branches are labelled by their monotonicity interval.  Complex branches
are labelled by a germ at a regular point and a homotopy class of continuation
path.  Near a ramification point of order \(m\), branches are indexed by
\(
 \omega_m^j=\exp(2\pi\ImaginaryUnit j/m),\ 0\le j<m
\).
Near infinity, a chosen branch of \(y^{1/m}\) is always stated explicitly.

\chapter{Ordinary inverse jets to arbitrary order}

\section{Lagrange--Bürmann in coefficient form}
Let
\[
 w=f(z)=c_1z+c_2z^2+c_3z^3+\cdots,
 \qquad c_1\ne0.
\]
Write \(z=g(w)=\sum_{n\ge1}b_nw^n\).  Lagrange--Bürmann gives
\begin{equation}
 b_n=\frac1n[z^{n-1}]
 \left(\frac{z}{f(z)}\right)^n.
 \label{eq:LB}
\end{equation}
The first coefficients are
\begin{align}
 b_1&=\frac1{c_1}, \\
 b_2&=-\frac{c_2}{c_1^3}, \\
 b_3&=\frac{2c_2^2}{c_1^5}-\frac{c_3}{c_1^4}, \label{eq:inverse-first3}\\
 b_4&=-\frac{5c_2^3}{c_1^7}
       +\frac{5c_2c_3}{c_1^6}-\frac{c_4}{c_1^5}.
\end{align}

\begin{proposition}[\statusproved\ Triangular reversion recurrence]
If \(b_1=1/c_1\), then for \(n\ge2\),
\[
 b_n=-\frac1{c_1}
 \sum_{k=2}^{n}c_k
 \sum_{\substack{j_1+\cdots+j_k=n\\j_i\ge1}}
 b_{j_1}\cdots b_{j_k}.
\]
This computes the inverse jet using only additions and multiplications in the
coefficient ring with \(c_1^{-1}\) adjoined.
\end{proposition}
\begin{proof}
Extract the coefficient of \(w^n\) from
\(f(g(w))=w\).  The contribution of \(c_1g\) is \(c_1b_n\); every other term
contains only \(b_j\) with \(j<n\).
\end{proof}

This recurrence is the default arbitrary-order mechanism throughout the
report.  It works equally well when each \(c_j\) is a rational function,
polylogarithm, Bernoulli polynomial, or formal series in auxiliary parameters.

\section{Bell-polynomial packaging}
If the forward jet is exponentiated,
\[
 \frac{F(p_0+z)}{F(p_0)}
 =\exp\left(\sum_{r\ge1}\ell_r\frac{z^r}{r!}\right),
\]
then its ordinary coefficients are complete exponential Bell polynomials:
\[
 [z^n]\frac{F(p_0+z)}{F(p_0)}
 =\frac1{n!}B_n(\ell_1,\ldots,\ell_n).
\]
Reversion can therefore be organized as the composition
\[
 \text{logarithmic derivatives}
 \longrightarrow \text{Bell polynomial}
 \longrightarrow \text{triangular inverse recurrence}.
\]
This is especially efficient for products because their logarithmic
derivatives are short power sums.

\section{Mixed parameter jets}
For \(F(p,\lambda)=y\), write
\(p=p_0+\sum_{r+s\ge1}b_{rs}u^rv^s\), where
\(u=y-y_0\) and \(v=\lambda-\lambda_0\).  Substitution into the bivariate Taylor
series of \(F\) gives a triangular recurrence ordered by total degree.  The
first mixed derivatives are
\begin{align}
 G_{yy}&=-\frac{F_{pp}}{F_p^3}, \\
 G_{y\lambda}&=
 -\frac{F_{p\lambda}}{F_p^2}
 +\frac{F_{pp}F_{\lambda}}{F_p^3}, \\
 G_{\lambda\lambda}&=-\frac{F_{\lambda\lambda}}{F_p}
 +2\frac{F_{p\lambda}F_{\lambda}}{F_p^2}
 -\frac{F_{pp}F_{\lambda}^2}{F_p^3}.
\end{align}
They quantify how an inverse in one parameter deforms when another parameter
moves, and they are indispensable near \(q=1/2\) in Fabius applications.

\chapter{Singular inverses: Puiseux, logarithmic, and multiscale regimes}

\section{Finite ramification}
Suppose
\[
 F(p_0+z)=y_0+c_mz^m+c_{m+1}z^{m+1}+\cdots,
 \qquad c_m\ne0,
\]
with \(m\ge2\).  Set
\(\xi=((y-y_0)/c_m)^{1/m}\).  Newton--Puiseux gives \(m\) branches
\begin{equation}
 p_j(y)=p_0+\omega_m^j\xi
 -\frac{c_{m+1}}{mc_m}\omega_m^{2j}\xi^2
 +\BigO(\xi^3),
 \qquad 0\le j<m.
 \label{eq:generic-puiseux}
\end{equation}
The expansion is convergent for algebraic functions and locally convergent for
holomorphic \(F\) after choosing a slit neighborhood.

\begin{proof}[Proof of the displayed coefficient]
Put \(z=\omega_m^j\xi+A\xi^2+\cdots\) and substitute into
\((F-y_0)/c_m=\xi^m\).  The coefficient of \(\xi^{m+1}\) is
\(m\omega_m^{j(m-1)}A+(c_{m+1}/c_m)\omega_m^{j(m+1)}\).  Since
\(\omega_m^{jm}=1\), solving gives the stated term.
\end{proof}

\section{Unfolding a fold}
If \(F_p=0\), \(F_{pp}\ne0\), and an auxiliary parameter \(\lambda\) varies,
then locally
\[
 y-y_*=A(p-p_*)^2+B(\lambda-\lambda_*)+
 C(p-p_*)(\lambda-\lambda_*)+\cdots.
\]
The discriminant curve in the \((y,\lambda)\)-plane is obtained by eliminating
\(p\) from \(F-y=0\) and \(F_p=0\).  Crossing it exchanges zero, one, or two
real branches.  This is the local model for the \(q=-1\) finite-product theorem
proved later.

\section{Essential and logarithmic inversion}
At \(q=1\), infinite products often have the form
\[
 \log F(\EulerE^{-\varepsilon})=-\frac{A}{\varepsilon}+B-C\varepsilon+
 D\varepsilon^3+\cdots.
\]
Let \(X=-\log y+B\).  Then
\begin{equation}
 \varepsilon=
 \frac{A}{X}+\frac{A^2C}{X^3}
 +\frac{A^4(-D)+2A^3C^2}{X^5}+\cdots,
 \label{eq:essential-reversion-template}
\end{equation}
with signs adjusted to the chosen forward convention.  Exponentiating gives
\(q=\EulerE^{-\varepsilon}\).  If \(\varepsilon\) also occurs inside a logarithm,
Lambert \(W\) or transseries rather than a pure inverse-power series appears.

\section{Newton polygons for multiple small scales}
If the target and an auxiliary parameter approach a singular point together,
a one-variable Taylor order is not intrinsic.  Write
\[
 F(p_0+z,\lambda_0+v)-y_0
 =\sum c_{rs}z^rv^s.
\]
A face of the Newton polygon determines a balance \(z\asymp v^\rho\).  Each
face polynomial supplies leading branches, after which a Puiseux recurrence
continues them.  The fixed-target inverse of \((a;q)_\infty\) near \(q=1\) is a
particularly important double-scaling example: the correct balance is
\(a\asymp\varepsilon\), not fixed \(a\).

\chapter{Simultaneous inversion and Good's multivariate formula}

\section{Square systems}
Let \(\bm F:\ComplexNumbers^r\to\ComplexNumbers^r\) and suppose
\(\det D\bm F(\bm p_0)\ne0\).  Then \(\bm p=\bm G(\bm y)\) exists locally and
\[
 D\bm G=(D\bm F)^{-1}.
\]
For example, recovering \((a,q)\) from two measurements
\[
 y_1=\QPochhammer{a}{q}{n},
 \qquad
 y_2=\QPochhammer{aq^s}{q}{m}
\]
is locally possible precisely when the two logarithmic sensitivity vectors are
not parallel.  This replaces the underdetermined scalar inverse by an
identifiable two-channel experiment.

\section{Good's multivariate Lagrange inversion}
Suppose formal series \(w_i\) satisfy
\[
 w_i=t_i\phi_i(\bm w),
 \qquad \phi_i(\bm0)\ne0,
 \qquad 1\le i\le r.
\]
Good's formula \cite{Good1960} states that, for a formal series \(H\),
\begin{align}
 H(\bm w(\bm t))
 =\sum_{\bm n\ge\bm0}\bm t^{\bm n}
 [\bm z^{\bm n}]\;H(\bm z)
 \prod_{i=1}^r\phi_i(\bm z)^{n_i}
 \det\left(\delta_{ij}-
 \frac{z_j}{\phi_i(\bm z)}\frac{\partial\phi_i}{\partial z_j}(\bm z)
 \right)_{i,j=1}^r.
 \label{eq:good}
\end{align}
It is the natural coefficient engine when the recovered base, argument, and
order are constrained by several q-product observations.  The determinant is
the multivariate Jacobian correction that has no scalar analogue.

\section{A practical jet algebra}
For computation, a multivariate jet is stored as a sparse map from multi-indices
to exact coefficients.  Composition is truncated by total or weighted degree;
Newton iteration doubles the known order:
\[
 \bm p_{new}=\bm p-(D\bm F(\bm p))^{-1}
 ( \bm F(\bm p)-\bm y).
\]
Over rational-function coefficient fields this gives exact jets.  Over
floating-point balls it gives certified jets provided the Jacobian inverse and
remainder bounds are enclosed.

\part{Finite \(q\)-Pochhammer products}

\chapter{Inversion in the argument}

\section{Forward algebra and logarithmic derivatives}
For \(n\in\NaturalNumbers\), set
\[
 P_n(a,q)=\QPochhammer{a}{q}{n}=\prod_{j=0}^{n-1}(1-aq^j),
 \qquad P_0=1.
\]
The finite \(q\)-binomial theorem gives
\begin{equation}
 P_n(a,q)=\sum_{m=0}^n(-1)^m q^{\binom m2}
 \GaussianBinomial{n}{m}{q} a^m.
 \label{eq:finite-p-in-a}
\end{equation}
Whenever no factor vanishes,
\begin{equation}
 \partial_a^r\log P_n(a,q)
 =-(r-1)!\sum_{j=0}^{n-1}
 \frac{q^{jr}}{(1-aq^j)^r},
 \qquad r\ge1.
 \label{eq:a-log-derivatives}
\end{equation}
Thus every regular argument-inverse coefficient is an explicit rational
function of finite geometric power sums.

\section{The branch issuing from \(a=0\)}
Let \(s=1-y\).  From \cref{eq:finite-p-in-a},
\[
 s=[n]_q a-q\GaussianBinomial{n}{2}{q}a^2+q^3\GaussianBinomial{n}{3}{q}a^3+\cdots,
\]
where \([n]_q=(1-q^n)/(1-q)\).  Applying
\cref{eq:inverse-first3} gives the following useful seed branch.

\begin{theorem}[\statusproved\ Argument inverse at the unit target]
If \([n]_q\ne0\), the branch satisfying \(a(1)=0\) is
\begin{align}
 a={}&\frac{s}{[n]_q}
 +\frac{q\GaussianBinomial{n}{2}{q}}{[n]_q^3}s^2 \notag \\
 &+\left(
 \frac{2q^2\GaussianBinomial{n}{2}{q}^2}{[n]_q^5}
 -\frac{q^3\GaussianBinomial{n}{3}{q}}{[n]_q^4}
 \right)s^3+\BigO(s^4).
 \label{eq:a-inverse-near-zero}
\end{align}
To arbitrary order,
\[
 [s^m]a(s)=\frac1m[z^{m-1}]
 \left(\frac{z}{1-P_n(z,q)}\right)^m.
\]
\end{theorem}
\begin{proof}
The displayed forward series has nonzero linear coefficient \([n]_q\).  Apply
Lagrange--Bürmann and then substitute the first three coefficients.
\end{proof}

The singular set \([n]_q=0\) is already informative: it consists of nontrivial
\(n\)-th roots of unity.  There the branch issuing from \(a=0\) ceases to be
ordinary; the first nonzero elementary symmetric coefficient determines its
Puiseux order.

\section{Inverse jets at the exact zeros}
For \(q\ne0\), the zeros in \(a\) are \(a_j=q^{-j}\),
\(0\le j<n\).  If they are distinct, each supplies an ordinary inverse branch
at target \(y=0\).

\begin{theorem}[\statusproved\ Exact derivative at a geometric zero]
Assume \(q^r\ne1\) for \(1\le r<n\).  Then
\begin{equation}
 P_{n,a}(q^{-j},q)=
 (-1)^{j+1}q^{-j(j-1)/2}(q;q)_j(q;q)_{n-1-j}.
 \label{eq:derivative-geometric-zero}
\end{equation}
Consequently the branch through \(a_j=q^{-j}\) begins
\begin{equation}
 a(y)=q^{-j}+
 \frac{(-1)^{j+1}q^{j(j-1)/2}}
 {(q;q)_j(q;q)_{n-1-j}}\,y+\BigO(y^2).
 \label{eq:inverse-at-geometric-zero}
\end{equation}
\end{theorem}
\begin{proof}
Factor the vanishing term:
\[
 P_{n,a}(q^{-j},q)=-q^j\prod_{\ell\ne j}(1-q^{\ell-j}).
\]
For \(\ell<j\), write
\(1-q^{\ell-j})=-q^{-(j-\ell)}(1-q^{j-\ell})\).  Their product is
\((-1)^jq^{-j(j+1)/2}(q;q)_j\); the remaining factors give
\((q;q)_{n-1-j}\).  Multiplication by \(-q^j\) proves
\cref{eq:derivative-geometric-zero}, and the inverse-function theorem gives
\cref{eq:inverse-at-geometric-zero}.
\end{proof}

The quadratic term follows without expanding the whole polynomial.  If
\(R_j(a)=\prod_{\ell\ne j}(1-aq^\ell)\), then
\[
 \frac{P_{n,aa}(a_j,q)}{P_{n,a}(a_j,q)}
 =2\frac{R_j'(a_j)}{R_j(a_j)}
 =-2\sum_{\ell\ne j}\frac{q^\ell}{1-q^{\ell-j}}.
\]
Substitute this into \(a=a_j+y/P_a-P_{aa}y^2/(2P_a^3)+\cdots\).

\section{Positive real base: complete real branch skeleton}
For \(0<q<1\), the zeros
\(1<q^{-(n-1)}<\cdots<q^{-1}<q^0=1\) are positive and simple.  Rolle's
theorem implies that every zero of \(\partial_aP_n\) is real, simple, and lies
strictly between consecutive zeros of \(P_n\).  Therefore the real inverse
covering is obtained by cutting the \(a\)-axis at these critical points.  The
critical values determine how many real preimages a target possesses.  This
recovers the branch skeleton used in the earlier inverse report, but
\cref{eq:derivative-geometric-zero} adds exact local normalizations at every
zero.

\section{Argument inverse at infinity}
Let \(N=\binom n2\).  The two leading powers of \cref{eq:finite-p-in-a} are
\[
 P_n(a,q)=(-1)^nq^Na^n
 +(-1)^{n-1}q^{(n-1)(n-2)/2}[n]_q a^{n-1}+\BigO(a^{n-2}).
\]
Hence, for a chosen \(n\)-th root
\(t=(y/((-1)^nq^N))^{1/n}\),
\begin{equation}
 a=t+\frac{[n]_q}{nq^{n-1}}+\BigO(t^{-1}).
 \label{eq:a-inverse-infinity}
\end{equation}
There are \(n\) branches, obtained by multiplying \(t\) by the \(n\)-th roots
of unity.  The complete descending series is obtained by substituting
\(a=t+\sum_{r\ge0}c_rt^{-r}\) into \cref{eq:finite-p-in-a}; its recurrence is
triangular.

\chapter{Base inversion near \(q=0\) and \(q=1\)}

\section{The small-base combinatorial jet}
Separate the \(j=0\) factor:
\[
 P_n(a,q)=(1-a)\prod_{j=1}^{n-1}(1-aq^j).
\]
The coefficient of \(q^r\) in the normalized product is the finite subset sum
\begin{equation}
 [q^r]\frac{P_n(a,q)}{1-a}
 =\sum_{\substack{S\subseteq\{1,\ldots,n-1\}\\sum_{j\in S}j=r}}
 (-a)^{\AbsoluteValue S}.
 \label{eq:subset-small-q}
\end{equation}
This supplies all forward and inverse coefficients without differentiation.

For \(n\ge2\) and \(a\notin\{0,1\}\), put
\[
 r=\frac{(1-a)-y}{a(1-a)}.
\]
The leading relation is \(r=q+\BigO(q^2)\), so the base inverse is ordinary.

\begin{proposition}[\statusproved\ Small-base inverse through cubic order]
For the branch \(q\to0\),
\[
\begin{array}{rll}
 n=2:&q=r,&\text{exactly},\\[0.25em]
 n=3:&q=r-r^2+(2+a)r^3+\BigO(r^4),&\\[0.25em]
 n\ge4:&q=r-r^2+(1+a)r^3+\BigO(r^4).&
\end{array}
\]
All later coefficients are obtained by applying the triangular reversion
recurrence to \cref{eq:subset-small-q}.
\end{proposition}
\begin{proof}
For \(n=3\),
\(r=q+q^2-aq^3+\BigO(q^4)\).  For \(n\ge4\), the subsets contributing through
degree three give
\(r=q+q^2+(1-a)q^3+\BigO(q^4)\).  Reversion proves the formulas.
\end{proof}

The excluded cases are structural, not removable bookkeeping.  At \(a=0\),
\(P_n\equiv1\); at \(a=1\), \(P_n\equiv0\) because the \(j=0\) factor vanishes.
For \(n=1\), the product is independent of \(q\).

\section{The logarithmic jet at \(q=1\)}
Set \(q=\EulerE^{-\varepsilon}\), and let
\(S_m(n)=\sum_{j=0}^{n-1}j^m\).  For \(|a|<1\), expansion of each logarithm and
interchange of finite sums gives
\begin{theorem}[\statusproved\ Complete finite-product jet at \(q=1\)]
\begin{equation}
 \log P_n(a,\EulerE^{-\varepsilon})
 =n\log(1-a)+\sum_{m\ge1}
 \frac{(-1)^{m+1}}{m!}
 \Polylogarithm_{1-m}(a)S_m(n)\varepsilon^m.
 \label{eq:finite-q1-logjet}
\end{equation}
The identity extends by analytic continuation in \(a\) on any simply connected
domain avoiding the moving zero set.
\end{theorem}
\begin{proof}
Use
\[
 \log(1-a\EulerE^{-j\varepsilon})
 =-\sum_{r\ge1}\frac{a^r}{r}\EulerE^{-jr\varepsilon}
\]
and expand the exponential.  The coefficient of \(\varepsilon^m\) is
\(-1)^{m+1}\sum_r a^rr^{m-1}S_m(n)/m!\), which is the displayed
polylogarithm.
\end{proof}

Let
\[
 w=\log\frac{y}{(1-a)^n}=c_1\varepsilon+c_2\varepsilon^2+c_3\varepsilon^3+\cdots,
 \qquad
 c_m=\frac{(-1)^{m+1}}{m!}\Polylogarithm_{1-m}(a)S_m(n).
\]
If \(c_1=S_1(n)a/(1-a)\ne0\), then
\begin{equation}
 \varepsilon=
 \frac{w}{c_1}-\frac{c_2}{c_1^3}w^2
 +\left(\frac{2c_2^2}{c_1^5}-\frac{c_3}{c_1^4}\right)w^3
 +\BigO(w^4),
 \qquad q=\EulerE^{-\varepsilon}.
 \label{eq:finite-q1-base-inverse}
\end{equation}
The formula is uniform on compact auxiliary-parameter sets on which \(c_1\)
stays away from zero.

\section{Inverting the argument while \(q\to1\)}
A different inverse problem keeps the target \(y\) fixed and recovers \(a\).
Let
\[
 a_0=1-y^{1/n},
\]
with a chosen branch of the \(n\)-th root.  Solving
\(P_n(a(\varepsilon),\EulerE^{-\varepsilon})=y\) using
\cref{eq:finite-q1-logjet} gives the following mixed jet.

\begin{theorem}[\statusproved\ Fixed-target argument inverse near \(q=1\)]
If \(a_0\ne1\), then
\begin{align}
 a(\varepsilon)={}&a_0+\frac{n-1}{2}a_0\varepsilon \notag \\
 &+\frac{a_0(n-1)
 \bigl(2n-4-3a_0(n-1)\bigr)}{24(1-a_0)}\varepsilon^2
 +\BigO(\varepsilon^3).
 \label{eq:finite-a-q1-mixed}
\end{align}
The higher coefficients are rational functions of \(a_0,n\) generated
triangularly from \cref{eq:finite-q1-logjet}.
\end{theorem}
\begin{proof}
Substitute \(a=a_0+A\varepsilon+B\varepsilon^2+\cdots\) into the logarithm of
the defining equation.  The constant term is \(n\log(1-a_0)=\log y\).  The
linear equation gives \(A=(n-1)a_0/2\); the quadratic equation gives the stated
\(B\).
\end{proof}

\chapter{The inverse atlas at \(q=-1\) and at roots of unity}

\section{Generic regularity at \(q=-1\)}
At the alternating base,
\begin{equation}
 P_n(a,-1)=(1-a)^{\lceil n/2\rceil}(1+a)^{\lfloor n/2\rfloor}.
 \label{eq:p-at-minus-one}
\end{equation}
Assume temporarily \(a\notin\{0,\pm1\}\).  Put
\[
 E=\sum_{\substack{1\le j<n\\j\text{ even}}}j,
 \qquad
 O=\sum_{\substack{1\le j<n\\j\text{ odd}}}j.
\]
Then
\begin{equation}
 \left.\partial_q\log P_n(a,q)\right|_{q=-1}
 =a\left(\frac{E}{1-a}-\frac{O}{1+a}\right).
 \label{eq:minus-one-log-first}
\end{equation}
Here
\(E=e(e+1)\), \(e=\lfloor(n-1)/2\rfloor\), and
\(O=o^2\), \(o=\lceil(n-1)/2\rceil\).  Unless the right side vanishes, the
inverse \(q(y)\) is an ordinary Taylor series based at
\(y_0=P_n(a,-1)\).

\section{The unique nontrivial critical parameter}
\begin{theorem}[\statusproved\ Complete finite-product criticality at
\(q=-1\)]
Let \(n\ge2\) and \(a\notin\{0,\pm1\}\).  Then
\(\partial_qP_n(a,-1)=0\) if and only if
\begin{equation}
 a=a_*=
 \begin{cases}
 \dfrac1{n-1},&n\text{ even},\\[0.45em]
 -\dfrac1n,&n\text{ odd}.
 \end{cases}
 \label{eq:a-star-minus-one}
\end{equation}
Write \(n=2s\) or \(n=2s+1\).  At this parameter,
\begin{align}
 \left.\partial_q^2\log P_{2s}(a_*,q)\right|_{q=-1}
 &=-\frac{(s+1)(2s-1)^2}{12s(s-1)}, \qquad s\ge2,
 \label{eq:L2-even} \\
 \left.\partial_q^2\log P_{2s+1}(a_*,q)\right|_{q=-1}
 &=-\frac{(s-1)(2s+1)^2}{12s(s+1)}, \qquad s\ge1.
 \label{eq:L2-odd}
\end{align}
Thus every case \(n\ge4\) is a strict quadratic fold.  The case \(n=3\) is
cubic and satisfies the exact identity
\begin{equation}
 P_3(-1/3,q)=\frac{32}{27}+\frac4{27}(q+1)^3.
 \label{eq:exact-cubic-minus-one}
\end{equation}
\end{theorem}
\begin{proof}
For \(a\ne0\), equation \cref{eq:minus-one-log-first} vanishes precisely when
\(E(1+a)=O(1-a)\), hence \(a=(O-E)/(O+E)\).  If \(n=2s\), then
\(E=s(s-1)\), \(O=s^2\), giving \(a=1/(2s-1)\).  If \(n=2s+1\), then
\(E=s(s+1)\), \(O=s^2\), giving \(a=-1/(2s+1)\).

Differentiate
\[
 \partial_q\log P_n=-a\sum_{j=1}^{n-1}
 \frac{jq^{j-1}}{1-aq^j}
\]
once more.  Separate even and odd \(j\) and use the polynomial sums of
\(j(j-1)\) and \(j^2\).  Substitution of \cref{eq:a-star-minus-one} simplifies
to \cref{eq:L2-even} and \cref{eq:L2-odd}.  The even expression is nonzero for
\(s\ge2\); the odd expression is nonzero for \(s\ge2\).  When \(s=1\), direct
multiplication gives \cref{eq:exact-cubic-minus-one}.
\end{proof}

\begin{corollary}[\statusproved\ Explicit local inverse types]
Let \(y_*=P_n(a_*,-1)\).  For \(n\ge4\), put \(L_2\) equal to the appropriate
quantity in \cref{eq:L2-even,eq:L2-odd}.  The two branches are
\begin{equation}
 q=-1\pm\sqrt{\frac{2(y-y_*)}{y_*L_2}}+\BigO(y-y_*).
 \label{eq:minus-one-fold-inverse}
\end{equation}
Since \(y_*>0\) and \(L_2<0\), the two nearby real branches occur for
\(y<y_*\).  For \(n=3\),
\begin{equation}
 q=-1+\omega_3^j\left(\frac{27}{4}
 \left(y-\frac{32}{27}\right)\right)^{1/3},
 \qquad j=0,1,2,
 \label{eq:minus-one-cubic-inverse}
\end{equation}
exactly.
\end{corollary}

\section{The singular endpoint parameters \(a=\pm1\)}
The parameter \(a=1\) makes \(P_n\equiv0\) because the \(j=0\) factor vanishes;
no scalar base inverse is identifiable.  For \(a=-1\), however, only the odd
factors vanish at \(q=-1\).  Let \(m=\lfloor n/2\rfloor\).  Then
\begin{theorem}[\statusproved\ Higher ramification of \((-1;q)_n\)]
\begin{equation}
 (-1;q)_n=
 2^{\lceil n/2\rceil}(2m-1)!!\,(q+1)^m
 \bigl(1+\BigO(q+1)\bigr).
 \label{eq:minus-one-a-minus-one}
\end{equation}
Consequently the inverse at target zero has \(m\) Puiseux branches,
\[
 q=-1+\omega_m^j
 \left(\frac{y}{2^{\lceil n/2\rceil}(2m-1)!!}\right)^{1/m}
 +\BigO(y^{2/m}).
\]
\end{theorem}
\begin{proof}
For odd \(j=2r-1\),
\(1+q^j=(2r-1)(q+1)+\BigO((q+1)^2)\).  Every even factor, including \(j=0\),
tends to \(2\).  Multiply the leading factors.
\end{proof}

\section{Arbitrary roots of unity}
Let \(\zeta\) be a root of unity and assume \(a\ne1\).  A factor with \(j>0\)
vanishes at \(q=\zeta\) exactly when \(a\zeta^j=1\).  Each such factor has a
simple zero in \(q\).  Hence
\begin{equation}
 \ord_{q=\zeta}P_n(a,q)
 =\#\{j\in\{1,\ldots,n-1\}:a\zeta^j=1\}.
 \label{eq:p-root-unity-multiplicity}
\end{equation}
If the count is \(m>0\), the inverse at target zero has \(m\) branches of the
form \cref{eq:generic-puiseux}.  If the count is zero but
\(\partial_qP_n(a,\zeta)=0\), the target is a nonzero critical value; its
ramification is found from the first nonzero derivative.  Eliminating \(a\)
from \(P_q=0\), \(P_{qq}=0\), and higher equations produces the root-of-unity
critical strata.

\chapter{Base inversion at \(q=\pm\infty\)}

\section{Reciprocal factorization}
For \(n\ge2\), \(a\notin\{0,1\}\), and \(N=\binom n2\), factor every term with
\(j\ge1\):
\begin{equation}
 P_n(a,q)=C(a)q^N
 \prod_{j=1}^{n-1}(1-a^{-1}q^{-j}),
 \qquad
 C(a)=(1-a)(-a)^{n-1}.
 \label{eq:p-reciprocal-factorization}
\end{equation}
Thus infinity is an ordinary point after the ramified coordinate
\(Y=(y/C)^{1/N}\) is introduced.

\begin{theorem}[\statusproved\ First large-base inverse terms]
On any chosen \(N\)-th-root branch,
\begin{equation}
 q=Y+\frac1{aN}+\BigO(Y^{-1}),
 \qquad
 Y=\left(\frac{y}{C(a)}\right)^{1/N}.
 \label{eq:p-large-q-inverse}
\end{equation}
The full expansion is a convergent Laurent series in \(Y^{-1}\) sufficiently
near infinity on the corresponding sheet.
\end{theorem}
\begin{proof}
The product in \cref{eq:p-reciprocal-factorization} equals
\(1-a^{-1}q^{-1}+\BigO(q^{-2})\).  Put \(q=Y+c+\BigO(Y^{-1})\).  Comparing the
coefficient of \(Y^{-1}\) in
\((q/Y)^N(1-a^{-1}q^{-1}+\cdots)=1\) gives \(Nc-a^{-1}=0\).
\end{proof}

For \(q\to-\infty\), the same formula applies with a branch of
\(Y=(y/C)^{1/N}\) tending to the negative real axis.  Whether a real branch is
available depends on the sign of \(C(-1)^N\) and on the target direction.  The
complex atlas is uniform: all \(N\) branches are related by roots of unity.

\section{Reciprocal transport as an algorithm}
Rather than expanding a high-degree polynomial at infinity directly, compute
the small-\(Q\) product
\[
 R(Q)=\prod_{j=1}^{n-1}(1-a^{-1}Q^j),
\]
revert \(Y=qR(1/q)^{1/N}\), and finally set \(q=Q^{-1}\).  This reuses the
small-base subset-sum machinery and prevents catastrophic cancellation.

\chapter{Order inversion, multiple arguments, and mixed recovery}

\section{Complex order}
For \(0<\AbsoluteValue q<1\), define
\begin{equation}
 (a;q)_\alpha=\frac{(a;q)_\infty}{(aq^\alpha;q)_\infty},
 \label{eq:complex-order}
\end{equation}
with a chosen branch of \(q^\alpha=\EulerE^{\alpha\PrincipalLogarithm q}\).  Its Lambert-series
form is
\begin{equation}
 \log(a;q)_\alpha
 =-\sum_{r\ge1}\frac{a^r(1-q^{\alpha r})}{r(1-q^r)}.
 \label{eq:complex-order-lambert}
\end{equation}

\begin{theorem}[\statusproved\ All order derivatives]
On every domain of absolute convergence,
\begin{equation}
 \partial_\alpha^m\log(a;q)_\alpha
 =(\PrincipalLogarithm q)^m\sum_{r\ge1}
 \frac{r^{m-1}a^rq^{\alpha r}}{1-q^r},
 \qquad m\ge1.
 \label{eq:alpha-derivatives}
\end{equation}
\end{theorem}
\begin{proof}
Differentiate \cref{eq:complex-order-lambert} term by term.
\end{proof}

If the first derivative is nonzero, every coefficient of the local inverse
\(\alpha(y)\) follows from \cref{eq:LB} and \cref{eq:alpha-derivatives}.  If it
vanishes, the first nonzero higher Lambert sum determines a Puiseux branch.
The branch also depends on the logarithm used in \(q^\alpha\); shifting
\(\PrincipalLogarithm q\) by \(2\pi\ImaginaryUnit\) changes the global order covering.

\section{Near \(q=1\): Bernoulli-polynomial order jet}
The generating identity
\begin{equation}
 \frac{1-\EulerE^{-\alpha x}}{1-\EulerE^{-x}}
 =\sum_{j\ge0}
 \frac{B_{j+1}(1)-B_{j+1}(1-\alpha)}{(j+1)!}x^j
 \label{eq:bernoulli-ratio}
\end{equation}
gives, for fixed \(|a|<1\),
\begin{equation}
 \log(a;\EulerE^{-\varepsilon})_\alpha
 =-\sum_{j\ge0}
 \frac{B_{j+1}(1)-B_{j+1}(1-\alpha)}{(j+1)!}
 \Polylogarithm_{1-j}(a)\varepsilon^j.
 \label{eq:complex-order-q1}
\end{equation}
The leading term is \(\alpha\log(1-a)\), as required.  If the target \(y\) is
fixed and
\[
 \alpha_0=\frac{\PrincipalLogarithm y}{\PrincipalLogarithm(1-a)},
\]
then, provided the chosen logarithms are consistent,
\begin{equation}
 \alpha(\varepsilon)=\alpha_0-
 \frac{\alpha_0(\alpha_0-1)\Polylogarithm_0(a)}{2\PrincipalLogarithm(1-a)}\varepsilon
 +\BigO(\varepsilon^2).
 \label{eq:alpha-inverse-q1}
\end{equation}
The coefficient is valid for complex \(\alpha_0\); no integrality is used.

\section{Discrete order and bracketing}
For integer order,
\[
 \frac{(a;q)_{n+1}}{(a;q)_n}=1-aq^n.
\]
On a real sign domain where these ratios stay in \((0,1)\) or in
\((1,\infty)\), the sequence is monotone and a discrete inverse order is
located by binary search.  The continuous inverse \cref{eq:complex-order}
then supplies a sub-integer interpolation and a Newton correction.  Outside
such domains, a discrete inverse is naturally set-valued.

\section{Several arguments and identifiability}
For
\[
 F(a_1,\ldots,a_s;q)=\prod_{i=1}^s(a_i;q)_{n_i},
\]
one scalar target leaves an \((s-1)\)-dimensional level set.  If all but \(a_r\)
are fixed, divide by the known factors and reuse the one-argument inverse.
If several arguments are to be recovered, use at least as many shifted-product
observations as unknowns.  The Jacobian entries are the finite sums in
\cref{eq:a-log-derivatives}; their determinant is an exact identifiability
criterion.

\part{The infinite product and complex-order limits}

\chapter{Argument inversion for \((a;q)_\infty\)}

\section{The principal real branch}
Assume \(0<q<1\) and define
\[
 P_\infty(a,q)=(a;q)_\infty=\prod_{j\ge0}(1-aq^j).
\]
For \(a<1\), all factors are positive and
\[
 \partial_a\log P_\infty(a,q)
 =-\sum_{j\ge0}\frac{q^j}{1-aq^j}<0.
\]
Moreover \(P_\infty(a,q)\to0\) as \(a\uparrow1\), while it tends to infinity
as \(a\to-\infty\).  Hence there is a unique real inverse
\[
 \Aq:(0,\infty)\longrightarrow(-\infty,1),
 \qquad
 (
 \Aq(y);q)_\infty=y.
\]
This is the principal inverse used in the earlier report.  The local formulas
below refine its endpoint normalization and extend directly to nonprincipal
complex branches.

\section{The branch at \(a=0\)}
Euler's expansion is
\begin{equation}
 (a;q)_\infty=\sum_{m\ge0}
 \frac{(-1)^mq^{\binom m2}}{(q;q)_m}a^m.
 \label{eq:euler-infinite-product}
\end{equation}
Let \(s=1-y\).  Reversion gives
\begin{align}
 a={}&(1-q)s+\frac{q(1-q)}{1+q}s^2 \notag \\
 &+(1-q)\left(
 \frac{2q^2}{(1+q)^2}
 -\frac{q^3}{(1+q)(1+q+q^2)}\right)s^3
 +\BigO(s^4).
 \label{eq:infinite-a-near-zero}
\end{align}
The arbitrary coefficient is again
\[
 [s^m]a(s)=\frac1m[z^{m-1}]
 \left(\frac{z}{1-(z;q)_\infty}\right)^m.
\]
The radius is limited by the nearest critical value of the entire function,
not merely by the nearest zero.

\section{All zero branches}
The zeros are \(a_m=q^{-m}\), \(m\ge0\), and are simple when \(q\) is not a
root of unity.  The finite-product calculation passes to the limit.

\begin{theorem}[\statusproved\ Exact derivative at every infinite-product zero]
For \(0<\AbsoluteValue q<1\) and \(m\ge0\),
\begin{equation}
 \left.\partial_a(a;q)_\infty\right|_{a=q^{-m}}
 =(-1)^{m+1}q^{-m(m-1)/2}(q;q)_m(q;q)_\infty.
 \label{eq:infinite-zero-derivative}
\end{equation}
Thus
\[
 a(y)=q^{-m}+
 \frac{(-1)^{m+1}q^{m(m-1)/2}}
 {(q;q)_m(q;q)_\infty}y+\BigO(y^2)
\]
on the branch through \(q^{-m}\).
\end{theorem}
\begin{proof}
Factor the zero at index \(m\).  The indices below \(m\) produce
\((-1)^mq^{-m(m+1)/2}(q;q)_m\), the indices above it produce
\((q;q)_\infty\), and differentiation of \(1-aq^m\) contributes \(-q^m\).
\end{proof}

\section{The principal branch at target zero}
Put \(a=1-\delta\).  Since
\[
 (a;q)_\infty=(1-a)(aq;q)_\infty,
\]
and \((q;q)_\infty\ne0\), the principal inverse has an ordinary endpoint jet.
Let
\[
 L_1(q)=\sum_{j\ge1}\frac{q^j}{1-q^j}.
\]
Then
\begin{equation}
 \Aq(y)=1-\frac{y}{(q;q)_\infty}
 +L_1(q)\left(\frac{y}{(q;q)_\infty}\right)^2
 +\BigO(y^3).
 \label{eq:principal-inverse-zero-target}
\end{equation}
Higher coefficients are polynomials in the Lambert sums
\(\sum_{j\ge1}q^{rj}/(1-q^j)^r\).

\section{Nonprincipal critical points}
Critical points solve
\[
 \sum_{j\ge0}\frac{q^j}{1-aq^j}=0.
\]
They lie between successive positive zeros on the real axis for \(0<q<1\),
and each simple critical point creates a square-root pair in the target plane.
The earlier inverse report develops their large-index asymptotics and
log-periodic organization.  The atlas viewpoint adds a practical rule: store
each critical germ together with the exact zero interval \((q^{-m},q^{-m-1})\)
and the continuation path, because ordering by target value can change at a
secondary-discriminant collision.

\chapter{Base inversion for the infinite product}

\section{At \(q=0\)}
For fixed \(a\notin\{0,1\}\), the first coefficients of the infinite product
coincide with those of every finite product of order at least four.  With
\[
 r=\frac{(1-a)-y}{a(1-a)},
\]
we therefore have
\begin{equation}
 q=r-r^2+(1+a)r^3+\BigO(r^4).
 \label{eq:infinite-base-zero-inverse}
\end{equation}
All orders follow from the subset-sum expansion over positive integers:
\[
 [q^N]\frac{(a;q)_\infty}{1-a}
 =\sum_{\substack{S\subset\PositiveIntegers\text{ finite}\\sum S=N}}(-a)^{\AbsoluteValue S}.
\]
This is a weighted distinct-partition polynomial in \(a\).

\section{Fixed argument and the essential singularity at \(q=1\)}
For fixed \(|a|<1\), set \(q=\EulerE^{-\varepsilon}\).  The Lambert series
\[
 \log(a;q)_\infty=-\sum_{r\ge1}\frac{a^r}{r(1-q^r)}
\]
and the Bernoulli expansion of \((1-\EulerE^{-x})^{-1}\) give
\begin{theorem}[\statusproved\ McIntosh-type logarithmic expansion]
\begin{equation}
 \log(a;\EulerE^{-\varepsilon})_\infty
 =-\frac{\Polylogarithm_2(a)}{\varepsilon}
 +\frac12\log(1-a)
 -\sum_{m\ge1}\frac{B_{2m}}{(2m)!}
 \Polylogarithm_{2-2m}(a)\varepsilon^{2m-1}.
 \label{eq:infinite-q1-forward}
\end{equation}
The equality is interpreted as an asymptotic expansion in sectors avoiding the
Stokes directions and the singularities of the chosen polylogarithm branches.
\end{theorem}
\begin{proof}
Insert
\[
 \frac1{1-\EulerE^{-x}}=\frac1x+\frac12+
 \sum_{m\ge1}\frac{B_{2m}}{(2m)!}x^{2m-1}
\]
into the Lambert series and sum powers of \(r\) as polylogarithms.  Standard
remainder estimates justify termwise truncation on compact subsets of
\(|a|<1\); analytic continuation gives the sectorial form.  This is the
classical asymptotic mechanism developed by McIntosh \cite{McIntosh1999}.
\end{proof}

For \(a\in(0,1)\), write
\[
 A=\Polylogarithm_2(a),
 \qquad B=\frac12\log(1-a),
 \qquad C=\frac1{12}\Polylogarithm_0(a),
 \qquad E=\frac1{720}\Polylogarithm_{-2}(a),
\]
and \(X=-\log y+B\).  Then \(y\to0^+\) corresponds to \(X\to\infty\), and
reversion of \cref{eq:infinite-q1-forward} gives
\begin{equation}
 \varepsilon=
 \frac{A}{X}+\frac{A^2C}{X^3}
 +\frac{2A^3C^2-A^4E}{X^5}
 +\BigO(X^{-7}),
 \qquad q=\EulerE^{-\varepsilon}.
 \label{eq:infinite-base-q1-inverse}
\end{equation}
This is the base inverse on the branch approaching \(q=1^-\).  Other complex
branches are obtained by the appropriate logarithmic and polylogarithmic
continuations.

\section{Monotonicity on principal sign domains}
If \(0<a<1\) and \(0<q<1\), then
\[
 \partial_q\log(a;q)_\infty
 =-a\sum_{j\ge1}\frac{jq^{j-1}}{1-aq^j}<0.
\]
Thus the product decreases from \(1-a\) at \(q=0\) to \(0\) at \(q=1^-\), and
for each \(0<y<1-a\) there is one positive real base.  If \(a<0\), the same
derivative is positive, so the product increases from \(1-a\) to infinity and
again has one positive real base in its range.  These global facts select the
principal branch to which \cref{eq:infinite-base-zero-inverse} and
\cref{eq:infinite-base-q1-inverse} belong.

\chapter{The fixed-target double-scaling inverse at \(q=1\)}

\section{Why fixed argument is the wrong scale}
For fixed \(a\in(0,1)\), \cref{eq:infinite-q1-forward} forces
\((a;q)_\infty\to0\) exponentially as \(q\to1^-\).  To keep a target
\(y\in(0,1)\) fixed, the argument must instead approach zero.  Balancing the
leading terms gives \(a\asymp\varepsilon\).  This regime is not obtained by
simply substituting into a fixed-\(a\) inverse after truncation; the infinitely
many powers of \(a\) and \(\varepsilon\) interact.

\section{Explicit expansion}
Let \(L=-\log y\), with \(L>0\), and seek
\[
 a(\varepsilon)=\sum_{m\ge1}A_m(L)\varepsilon^m.
\]

\begin{theorem}[\statusproved\ Five-term fixed-target inverse]
The principal solution of
\[
 (a;\EulerE^{-\varepsilon})_\infty=\EulerE^{-L}
\]
has the expansion
\begin{align}
 a={}&L\varepsilon-
 \frac{L(L+2)}4\varepsilon^2
 +\frac{L(L^2+9L+12)}{72}\varepsilon^3 \notag \\
 &-\frac{L(L^3+4L^2+12L+24)}{576}\varepsilon^4 \notag \\
 &-\frac{L(31L^4-75L^3-800L^2-720)}{86400}\varepsilon^5
 +\BigO(\varepsilon^6).
 \label{eq:double-scaling-five}
\end{align}
Every \(A_m(L)\) is a polynomial in \(L\) with rational coefficients and a
factor \(L\).
\end{theorem}
\begin{proof}
Start from the exact Lambert series
\[
 \log(a;\EulerE^{-\varepsilon})_\infty
 =-\sum_{r\ge1}\frac{a^r}{r(1-\EulerE^{-r\varepsilon})}.
\]
Insert the Bernoulli expansion of the denominator and the ansatz for \(a\).
At coefficient \(\varepsilon^{m-1}\), the new unknown \(A_m\) occurs only in
the term \(-a/\varepsilon\), with coefficient \(-1\).  The coefficient
equation is therefore triangular.  Solving successively gives the displayed
polynomials.  The same argument proves rational polynomiality and the factor
\(L\).
\end{proof}

\begin{algorithm}[Arbitrary-order double-scaling coefficients]
Fix an order \(M\).  Truncate
\[
 -\sum_{r=1}^{M+1}\frac{a^r}{r}
 \left(\frac1{r\varepsilon}+\frac12+
 \sum_{1\le j\le M/2}\frac{B_{2j}}{(2j)!}
 (r\varepsilon)^{2j-1}\right)
\]
at \(\varepsilon^M\), substitute
\(a=\sum_{m=1}^{M+1}A_m\varepsilon^m\), set the constant term to \(-L\), and
set every positive-power coefficient through \(\varepsilon^M\) to zero.  Solve
in the order \(A_1,A_2,\ldots\).  Only \(r\le M+1\) can contribute, so the
computation is finite and exact.
\end{algorithm}

\section{Complex targets and logarithm branches}
For a complex target \(y\ne0\), choose \(L=-\PrincipalLogarithm y\).  Each logarithm branch
produces a formal double-scaling branch.  The condition \(a\to0\) selects the
same local balance, but different values \(L+2\pi\ImaginaryUnit k\) lead to different
complex sheets.  Coalescence occurs at \(L=0\), where the fixed target is
\(y=1\) and the leading coefficient vanishes; a new Newton polygon is then
required.

\chapter{Complex order near \(q=0\), resonances, and loss of identifiability}

\section{A generalized power series rather than a Taylor series}
For \((a;q)_\alpha\) as in \cref{eq:complex-order}, suppose
\(\Re\alpha>0\).  From \cref{eq:complex-order-lambert},
\[
 \frac{1-q^{\alpha r}}{1-q^r}
 =(1-q^{\alpha r})(1+q^r+q^{2r}+\cdots).
\]
Unless \(\alpha\in\NaturalNumbers\), the expansion at \(q=0\) belongs to the generalized
power semigroup generated by \(1\) and \(\alpha\), not to \(\ComplexNumbers[[q]]\).  If
\(0<\Re\alpha<1\), the first order-dependent correction is \(q^\alpha\); if
\(\Re\alpha>1\), the first correction is \(q\) and is independent of
\(\alpha\), while order dependence enters later.

\begin{remark}[Resonance]
When \(\alpha=r/s\) is rational, the generalized series becomes an ordinary
Puiseux series in \(q^{1/s}\).  When \(\alpha\) is irrational, its support is a
nonlattice additive semigroup.  If two exponents coincide, coefficient
resonance changes the inverse recurrence.  Thus ``the expansion at \(q=0\)''
for complex order cannot be stated without the arithmetic type of
\(\alpha\).
\end{remark}

\section{Order identifiability collapses at zero base}
For every fixed \(\Re\alpha>0\),
\[
 \lim_{q\to0}(a;q)_\alpha=1-a.
\]
Therefore the Jacobian of the map \(\alpha\mapsto(a;q)_\alpha\) tends to zero.
Recovering \(\alpha\) from a noisy product value becomes singular as
\(q\to0\), even though the complex-order function itself has a finite limit.
This is a parameter-estimation obstruction, not a failure of analytic
continuation.

\section{Near integer order}
Put \(\alpha=n+\eta\).  From \cref{eq:alpha-derivatives},
\begin{align}
 \log(a;q)_{n+\eta}
 ={}&\log(a;q)_n
 +\eta(\PrincipalLogarithm q)\sum_{r\ge1}
 \frac{a^rq^{nr}}{1-q^r} \notag \\
 &+\frac{\eta^2}{2}(\PrincipalLogarithm q)^2\sum_{r\ge1}
 \frac{ra^rq^{nr}}{1-q^r}+\BigO(\eta^3).
 \label{eq:order-near-integer}
\end{align}
This supplies a series in the order parameter itself and a local inverse
\(\eta(y)\).  It is preferable to finite differencing integer products because
it retains exact branch information through \(\PrincipalLogarithm q\).

\part{Gaussian coefficients and palindromic \(q\)-polynomials}

\chapter{Gaussian base inversion near \(q=0\) and \(q=1\)}

\section{Restricted partitions at zero base}
For \(0<k<n\), let
\[
 G_{n,k}(q)=\GaussianBinomial{n}{k}{q},
 \qquad b=n-k,
 \qquad d=kb.
\]
Its coefficients count partitions fitting in a \(k\times b\) rectangle:
\begin{equation}
 G_{n,k}(q)=\sum_{r=0}^{d}p_{k,b}(r)q^r.
 \label{eq:gaussian-partitions}
\end{equation}
In particular \(p_{k,b}(0)=p_{k,b}(1)=1\), so the positive-base inverse at
target \(1\) is always ordinary.

Let \(u=y-1\), and set
\begin{align}
 c_2&=\ind_{k\ge2}+\ind_{b\ge2}, \\
 c_3&=\ind_{b\ge3}+\ind_{k\ge2,\,b\ge2}+\ind_{k\ge3}.
 \label{eq:c2c3-rectangle}
\end{align}
These are the numbers of rectangle-compatible partitions of \(2\) and \(3\).

\begin{theorem}[\statusproved\ Gaussian inverse at \(q=0\)]
For every \(0<k<n\),
\begin{equation}
 q=u-c_2u^2+(2c_2^2-c_3)u^3+\BigO(u^4).
 \label{eq:gaussian-q0-inverse}
\end{equation}
More generally,
\[
 [u^m]q(u)=\frac1m[z^{m-1}]
 \left(\frac{z}{G_{n,k}(z)-1}\right)^m.
\]
\end{theorem}
\begin{proof}
Equation \cref{eq:gaussian-partitions} begins
\(u=q+c_2q^2+c_3q^3+\cdots\).  Apply ordinary reversion.
\end{proof}

For a thin rectangle, say \(k=1\), this reduces to inversion of
\(1+q+\cdots+q^{n-1}\).  For \(k,b\ge3\), one has \(c_2=2\), \(c_3=3\), so
\(q=u-2u^2+5u^3+\cdots\).

\section{Bernoulli cumulants at unit base}
Using
\[
 G_{n,k}(q)=\prod_{j=1}^k\frac{1-q^{b+j}}{1-q^j}
\]
and \(q=\EulerE^{-\varepsilon}\), define
\[
 \Delta_m=S_m(n)-S_m(k)-S_m(b),
 \qquad S_m(N)=\sum_{j=1}^{N}j^m.
\]
The elementary expansion
\[
 \log\frac{1-\EulerE^{-x}}x
 =-\frac x2+\sum_{r\ge1}
 \frac{B_{2r}}{2r(2r)!}x^{2r}
\]
yields
\begin{theorem}[\statusproved\ Complete Gaussian logarithmic jet]
\begin{equation}
 \log G_{n,k}(\EulerE^{-\varepsilon})
 =\log\binom nk-\frac d2\varepsilon
 +\sum_{r\ge1}
 \frac{B_{2r}\Delta_{2r}}{2r(2r)!}\varepsilon^{2r}.
 \label{eq:gaussian-q1-forward}
\end{equation}
\end{theorem}
\begin{proof}
Apply the displayed scalar expansion to each numerator and denominator factor.
The constant terms give \(\binom nk\); the linear terms telescope to
\(-kb\varepsilon/2\); the even power sums give \(\Delta_{2r}\).
\end{proof}

Put \(v=\log(y/\binom nk)\).  Since the linear coefficient is nonzero,
ordinary reversion gives an especially simple first jet.

\begin{corollary}[\statusproved\ Gaussian base inverse at \(q=1\)]
For \(0<k<n\),
\begin{equation}
 \varepsilon=-\frac{2v}{d}
 +\frac{\Delta_2}{3d^3}v^2
 -\frac{\Delta_2^2}{9d^5}v^3
 +\BigO(v^4),
 \qquad q=\EulerE^{-\varepsilon}.
 \label{eq:gaussian-q1-inverse}
\end{equation}
The coefficient of \(v^4\) is the first one involving \(\Delta_4\).
\end{corollary}
\begin{proof}
In \cref{eq:gaussian-q1-forward},
\(c_1=-d/2\), \(c_2=\Delta_2/24\), and \(c_3=0\).  Substitute these values in
\cref{eq:inverse-first3}.
\end{proof}

Since \(\Delta_2=d(n+1)\), the variance scale that controls the centered inverse
below is already visible in the quadratic coefficient.

\chapter{The centered-palindromic inverse theorem}

\section{A universal probabilistic normalization}
Let
\[
 P(q)=\sum_{j=0}^D c_jq^j,
 \qquad c_j\ge0,
 \qquad c_j=c_{D-j},
\]
with \(P\) nonconstant.  Define a random variable \(X\) by
\(\Pr(X=j)=c_j/P(1)\).  Palindromicity is exactly the distributional symmetry
\(X\overset d=D-X\).

For real \(\varepsilon\), set
\begin{equation}
 H_P(\varepsilon)=
 \frac{\EulerE^{D\varepsilon/2}P(\EulerE^{-\varepsilon})}{P(1)}
 =\Expectation\EulerE^{-\varepsilon(X-D/2)}.
 \label{eq:centered-palindromic}
\end{equation}

\begin{theorem}[\statusproved\ Global two-branch inverse]
The function \(H_P\) is even, satisfies \(H_P(0)=1\), is strictly increasing on
\((0,\infty)\), and tends to infinity there.  Hence for every \(h>1\) the
equation \(H_P(\varepsilon)=h\) has exactly two real solutions
\(\varepsilon=\pm E_P(h)\).  The inverse has a square-root branch point at
\(h=1\).
\end{theorem}
\begin{proof}
Let \(Z=X-D/2\).  Symmetry gives \(Z\overset d=-Z\), hence
\[
 H_P(\varepsilon)=\Expectation\cosh(\varepsilon Z).
\]
It is even and at least one.  For \(\varepsilon>0\),
\[
 H_P'(\varepsilon)=\Expectation[Z\sinh(\varepsilon Z)]>0,
\]
because \(z\sinh(\varepsilon z)\ge0\) and the nonconstant polynomial makes it
positive with nonzero probability.  Since \(c_0=c_D>0\), the endpoint terms
force exponential growth.  Finally,
\(H_P(\varepsilon)=1+\Variance(X)\varepsilon^2/2+\cdots\), so the local inverse is
square-root.
\end{proof}

This theorem resolves a common apparent paradox.  The unnormalized Gaussian
polynomial has an ordinary inverse at \(q=1\) because of its deterministic
linear drift \(D/2\).  Removing that drift exposes the reciprocal symmetry and
turns \(q=1\) into a genuine two-sheeted branch point.

\section{Cumulant Puiseux series}
Let \(\kappa_{2r}\) be the even cumulants of \(Z\); odd cumulants vanish.  Put
\(u=\log H_P(\varepsilon)\).  Then
\[
 u=\frac{\kappa_2}{2!}\varepsilon^2+
 \frac{\kappa_4}{4!}\varepsilon^4+
 \frac{\kappa_6}{6!}\varepsilon^6+\cdots.
\]
Reversion in \(\varepsilon^2\) gives
\begin{theorem}[\statusproved\ Centered inverse through order \(u^{5/2}\)]
\begin{align}
 \varepsilon={}&\pm\sqrt{\frac{2u}{\kappa_2}}
 \Biggl[
 1-\frac{\kappa_4}{12\kappa_2^2}u \notag \\
 &\qquad+
 \left(\frac{7\kappa_4^2}{288\kappa_2^4}
       -\frac{\kappa_6}{180\kappa_2^3}\right)u^2
 +\BigO(u^3)
 \Biggr].
 \label{eq:centered-puiseux}
\end{align}
\end{theorem}
\begin{proof}
Write \(u=A\varepsilon^2+B\varepsilon^4+C\varepsilon^6+\cdots\) and substitute
\(\varepsilon=(u/A)^{1/2}(1+\alpha u+\beta u^2+\cdots)\).  Solving the
coefficients gives
\(\alpha=-B/(2A^2)\) and
\(\beta=(-4AC+7B^2)/(8A^4)\).  Insert
\(A=\kappa_2/2\), \(B=\kappa_4/24\), \(C=\kappa_6/720\).
\end{proof}

\section{Gaussian cumulants}
For \(P=G_{n,k}\), comparison with \cref{eq:gaussian-q1-forward} gives
\begin{equation}
 \kappa_{2r}=\frac{B_{2r}}{2r}\Delta_{2r},
 \qquad
 \kappa_2=\frac{k(n-k)(n+1)}{12}.
 \label{eq:gaussian-cumulants}
\end{equation}
Thus \cref{eq:centered-puiseux} is an explicit Bernoulli--Faulhaber expansion.
The same construction applies to any positive palindromic q-enumerator; no
product representation is required.

\chapter{Roots of unity and the square-free Gaussian covering}

\section{Cyclotomic factorization}
Let \(\Phi_m(q)\) be the \(m\)-th cyclotomic polynomial.  Counting multiples of
\(m\) in the product formula yields
\begin{equation}
 G_{n,k}(q)=\prod_{m=2}^{n}\Phi_m(q)^{e_m},
 \qquad
 e_m=\left\lfloor\frac nm\right\rfloor
 -\left\lfloor\frac km\right\rfloor
 -\left\lfloor\frac{n-k}m\right\rfloor.
 \label{eq:gaussian-cyclotomic}
\end{equation}

\begin{theorem}[\statusproved\ Gaussian polynomials are square-free]
For \(0<k<n\), every exponent \(e_m\) in
\cref{eq:gaussian-cyclotomic} is either \(0\) or \(1\).  Consequently every
complex zero of \(G_{n,k}\) is a simple root of unity.
\end{theorem}
\begin{proof}
Write \(n=k+b\).  Then
\[
 e_m=\left\lfloor\frac{k+b}m\right\rfloor
 -\left\lfloor\frac km\right\rfloor
 -\left\lfloor\frac bm\right\rfloor.
\]
The difference records whether the two residues modulo \(m\) produce a carry,
and hence is \(0\) or \(1\).  Cyclotomic polynomials are separable in
characteristic zero and mutually coprime.
\end{proof}

If \(\zeta\) is primitive of order \(m\) and \(e_m=1\), write
\[
 R_m(q)=\prod_{r\ne m}\Phi_r(q)^{e_r}.
\]
Then
\begin{equation}
 G_{n,k}'(\zeta)=\Phi_m'(\zeta)R_m(\zeta)\ne0,
 \qquad
 q(y)=\zeta+\frac{y}{\Phi_m'(\zeta)R_m(\zeta)}+\BigO(y^2).
 \label{eq:gaussian-root-inverse}
\end{equation}
Thus Gaussian root-of-unity inverses are ordinary.  Higher Puiseux behavior at
roots of unity belongs to q-multinomials and q-factorials, not to Gaussian
polynomials themselves.

\section{The alternating base}
At \(q=-1\),
\begin{equation}
 G_{n,k}(-1)=
 \begin{cases}
 0,&n\text{ even and }k\text{ odd},\\[0.25em]
 \displaystyle\binom{\lfloor n/2\rfloor}{\lfloor k/2\rfloor},&\text{otherwise}.
 \end{cases}
 \label{eq:gaussian-minus-one-value}
\end{equation}
The zero case is always simple because \(e_2=1\).

\begin{theorem}[\statusproved\ Exact derivative of the alternating zero]
If \(n=2m\) and \(k=2r+1\), with \(0\le r<m\), then
\begin{equation}
 G_{2m,2r+1}'(-1)=(m-r)\binom mr.
 \label{eq:gaussian-minus-one-derivative}
\end{equation}
Hence
\[
 q=-1+\frac{y}{(m-r)\binom mr}+\BigO(y^2)
\]
for the inverse branch through the alternating zero.
\end{theorem}
\begin{proof}
In the product
\[
 G_{2m,2r+1}(q)=\prod_{j=1}^{2r+1}
 \frac{1-q^{2(m-r)-1+j}}{1-q^j},
\]
there is one more even numerator exponent than even denominator exponent.
For an even exponent \(2s\),
\((1-q^{2s})/(1+q)\to2s\); every odd factor tends to \(2\).  The ratio of the
even leading coefficients is
\(2m!/(r!(m-r-1)!)\), while the odd factors contribute \(1/2\).  Their product
is \((m-r)\binom mr\).
\end{proof}

\section{Nonzero targets at a root of unity}
Even when \(G(\zeta)\ne0\), the derivative may vanish and create a fold at a
nonzero target.  These points are zeros of
\(G'(q)\) constrained to the root-of-unity locus.  Their elimination ideal is
\[
 \langle \Phi_m(q),\;G'(q)\rangle.
\]
For fixed \((n,k,m)\) it is decidable by exact resultants.  A systematic table
of such critical coincidences is a natural continuation of the earlier
secondary-discriminant computations.

\chapter{Large base and reciprocal Gaussian branches}

\section{Reciprocity}
The rectangle complement involution gives
\begin{equation}
 G_{n,k}(q)=q^dG_{n,k}(q^{-1}),
 \qquad d=k(n-k).
 \label{eq:gaussian-reciprocity}
\end{equation}
Let
\[
 G_{n,k}(Q)=1+Q+c_2Q^2+c_3Q^3+\cdots
\]
with \(c_2,c_3\) as in \cref{eq:c2c3-rectangle}, and put
\(Y=y^{1/d}\) on a chosen branch.

\begin{theorem}[\statusproved\ Large-target Gaussian inverse]
\begin{align}
 q={}&Y-\frac1d
 +\left(\frac{d-1}{2d^2}-\frac{c_2}{d}\right)Y^{-1}
 +\BigO(Y^{-2}).
 \label{eq:gaussian-large-q}
\end{align}
The coefficient of \(Y^{-2}\) is
\[
 -\frac{(d-1)(2d-1)}{3d^3}
 +\frac{2c_2(d-1)}{d^2}-\frac{c_3}{d}.
\]
\end{theorem}
\begin{proof}
Substitute
\(q=Y(1+A_1Y^{-1}+A_2Y^{-2}+\cdots)\) into
\(q^dG(q^{-1})=y\) and compare powers of \(Y^{-1}\).  The equations are
triangular and give the displayed coefficients.
\end{proof}

For \(q\to-\infty\), reciprocity remains valid.  If \(d\) is odd, the leading
real target tends to \(-\infty\); if \(d\) is even, it tends to
\(+\infty\).  The complex branches are indexed by the \(d\)-th roots in \(Y\).

\chapter{Continuous upper and lower parameters; q-multinomials}

\section{q-gamma interpolation of the Gaussian coefficient}
For \(0<q<1\), define
\begin{equation}
 \mathcal G_q(\alpha,\beta)=
 \frac{\Gamma_q(\alpha+1)}
 {\Gamma_q(\beta+1)\Gamma_q(\alpha-\beta+1)}.
 \label{eq:continuous-gaussian}
\end{equation}
It agrees with \(\GaussianBinomial{n}{k}{q}\) at integers.  Let \(\psi_q\) denote the
\(q\)-digamma function.  Then
\begin{align}
 \partial_\alpha\log\mathcal G_q
 &=\psi_q(\alpha+1)-\psi_q(\alpha-\beta+1), \label{eq:G-alpha-derivative}\\
 \partial_\beta\log\mathcal G_q
 &=-\psi_q(\beta+1)+\psi_q(\alpha-\beta+1).
 \label{eq:G-beta-derivative}
\end{align}

\begin{theorem}[\statusproved\ Parameter branch geometry]
Assume \(\alpha>\beta>0\).
\begin{enumerate}
\item Since \(\psi_q\) is strictly increasing, \(\mathcal G_q\) is strictly
increasing in \(\alpha\).  Its upper-parameter inverse is therefore unique on
every interval contained in this domain.
\item As a function of \(\beta\in(0,\alpha)\), it is symmetric about
\(\alpha/2\), strictly increasing before the midpoint and strictly decreasing
after it.  At the midpoint,
\[
 \partial_\beta^2\log\mathcal G_q
 =-2\psi_q'(\alpha/2+1)<0.
\]
Thus the lower-parameter inverse has two real branches below the maximum and a
square-root merger at the maximum.
\end{enumerate}
\end{theorem}
\begin{proof}
For \(0<q<1\), \(\psi_q'(x)>0\).  The signs follow directly from
\cref{eq:G-alpha-derivative,eq:G-beta-derivative}; differentiating the second
identity proves the midpoint statement.
\end{proof}

Near \(q=1\), the q-gamma expansion derived later implies
\begin{equation}
 \log\mathcal G_{\EulerE^{-\varepsilon}}(\alpha,\beta)
 =\log\frac{\Gamma(\alpha+1)}
 {\Gamma(\beta+1)\Gamma(\alpha-\beta+1)}
 -\frac{\beta(\alpha-\beta)}2\varepsilon+\BigO(\varepsilon^2).
 \label{eq:continuous-gaussian-q1}
\end{equation}
Hence the base inverse is ordinary whenever
\(\beta(\alpha-\beta)\ne0\), and mixed inverses in \((q,\alpha,\beta)\) follow
from the Jacobian formulas of Part I.

\section{q-multinomial cyclotomic carry theorem}
For nonnegative integers \(k_1,\ldots,k_r\) with \(n=\sum_i k_i\), write
\[
 \qmulti{n}{k_1,\ldots,k_r}
 =\frac{(q;q)_n}{\prod_{i=1}^r(q;q)_{k_i}}.
\]
At a primitive \(d\)-th root, the cyclotomic exponent is
\begin{equation}
 e_d=\left\lfloor\frac nd\right\rfloor
 -\sum_{i=1}^r\left\lfloor\frac{k_i}d\right\rfloor.
 \label{eq:qmulti-cyclotomic}
\end{equation}

\begin{theorem}[\statusproved\ Residue-carry ramification]
Let \(r_i\equiv k_i\pmod d\), \(0\le r_i<d\).  Then
\begin{equation}
 e_d=\left\lfloor\frac{r_1+\cdots+r_r}{d}\right\rfloor,
 \qquad 0\le e_d\le r-1.
 \label{eq:carry-count}
\end{equation}
If \(e_d=m>0\), the q-multinomial has a zero of order \(m\) at every primitive
\(d\)-th root and its base inverse at target zero has \(m\) Puiseux branches.
\end{theorem}
\begin{proof}
Write \(k_i=d\ell_i+r_i\).  Since
\(n=d\sum_i\ell_i+\sum_i r_i\), substitution in
\cref{eq:qmulti-cyclotomic} gives \cref{eq:carry-count}.  The cyclotomic factor
\(\Phi_d^{e_d}\) has that exact zero order, and Newton--Puiseux gives the inverse
branches.
\end{proof}

For \(r=2\), the carry count is at most one, recovering Gaussian
square-freeness.  For many parts it can be arbitrarily large.  This is the
simplest exact mechanism by which passing from q-binomials to another standard
q-analog changes an ordinary root-of-unity inverse into a higher Puiseux
inverse.

\section{Centered q-multinomial inverse}
The q-multinomial has nonnegative palindromic coefficients and degree
\[
 D=\sum_{1\le i<j\le r}k_ik_j.
\]
Therefore \cref{eq:centered-palindromic} applies verbatim.  Its centered base
inverse has exactly two global real branches, and its local square-root
coefficients are the cumulants of the multiset-inversion statistic.  In
contrast, its root-of-unity inverse can have arbitrarily high finite
ramification by \cref{eq:carry-count}; the same polynomial thus exhibits two
different singular mechanisms at different targets.

\part{Other principal \(q\)-analogs}

\chapter{q-numbers and q-factorials}

\section{The q-number: one exact inverse and one algebraic inverse}
For \(q\ne1\), define
\[
 [x]_q=\frac{1-q^x}{1-q},
 \qquad q^x=\EulerE^{x\PrincipalLogarithm q}.
\]
Inversion in \(x\) is elementary once a logarithm branch is chosen:
\begin{equation}
 x=\frac{\PrincipalLogarithm\bigl(1-(1-q)y\bigr)}{\PrincipalLogarithm q}.
 \label{eq:qnumber-x-inverse}
\end{equation}
The branch points are inherited from both logarithms.  When \(x=m\in\NaturalNumbers\),
inversion in \(q\) is algebraic because
\[
 [m]_q=1+q+\cdots+q^{m-1}
\].

At \(q=0\), let \(u=y-1\).  Then
\[
\begin{array}{rll}
 m=2:&q=u,&\text{exactly},\\
 m=3:&q=u-u^2+2u^3+\BigO(u^4),&\\
 m\ge4:&q=u-u^2+u^3+\BigO(u^4).&
\end{array}
\]
At \(q=\EulerE^{-\varepsilon}\to1\),
\begin{align}
 [x]_{\EulerE^{-\varepsilon}}={}&x+
 \frac{x(1-x)}2\varepsilon
 +\frac{x(2x^2-3x+1)}{12}\varepsilon^2 \notag \\
 &-\frac{x^2(x-1)^2}{24}\varepsilon^3+\BigO(\varepsilon^4).
 \label{eq:qnumber-q1}
\end{align}
Thus, for \(x\notin\{0,1\}\),
\begin{equation}
 \varepsilon=
 \frac{2(y-x)}{x(1-x)}
 +\frac{2(2x-1)}{3x^2(x-1)^2}(y-x)^2
 +\BigO((y-x)^3).
 \label{eq:qnumber-q1-inverse}
\end{equation}
The values \(x=0\) and \(x=1\) are unidentifiable because the q-number is
identically \(0\) or \(1\).

For integer \(m\ge2\), the alternating base is regular:
\[
 [2r]_{-1}=0,
 \quad \left.\partial_q[2r]_q\right|_{q=-1}=r;
 \qquad
 [2r+1]_{-1}=1,
 \quad \left.\partial_q[2r+1]_q\right|_{q=-1}=-r.
\]
At infinity, \([m]_q=q^{m-1}(1+q^{-1}+\cdots)\), so the inverse has
\(m-1\) Puiseux branches.

\section{q-factorials at the distinguished bases}
Let
\[
 [n]_q!=\prod_{j=1}^n[j]_q=\Gamma_q(n+1).
\]
At zero base,
\begin{equation}
 [n]_q!=1+(n-1)q+\frac{(n-2)(n+1)}2q^2+\BigO(q^3).
 \label{eq:qfactorial-zero}
\end{equation}
For \(n\ge2\) and \(u=y-1\),
\begin{equation}
 q=\frac{u}{n-1}
 -\frac{(n-2)(n+1)}{2(n-1)^3}u^2+\BigO(u^3).
 \label{eq:qfactorial-zero-inverse}
\end{equation}

\begin{theorem}[\statusproved\ Exact alternating ramification of q-factorials]
Let \(m=\lfloor n/2\rfloor\).  Then
\begin{equation}
 [n]_q!=m!\,(q+1)^m\bigl(1+\BigO(q+1)\bigr)
 \qquad(q\to-1).
 \label{eq:qfactorial-minus-one}
\end{equation}
Hence the inverse at target zero is
\begin{equation}
 q=-1+\omega_m^j\left(\frac{y}{m!}\right)^{1/m}
 +\BigO(y^{2/m}),
 \qquad 0\le j<m.
 \label{eq:qfactorial-minus-one-inverse}
\end{equation}
\end{theorem}
\begin{proof}
Exactly the even q-numbers vanish.  For \(j=2r\),
\([2r]_q=r(q+1)+\BigO((q+1)^2)\), while every odd q-number tends to \(1\).
Multiplying \(r=1,\ldots,m\) proves the result.
\end{proof}

More generally, at a primitive \(d\)-th root \(\zeta\),
\begin{equation}
 \ord_{q=\zeta}[n]_q!=\left\lfloor\frac nd\right\rfloor,
 \qquad d>1.
 \label{eq:qfactorial-root-order}
\end{equation}
This follows because \([j]_q\) has a simple zero exactly when \(d\mid j\).
Thus q-factorials supply arbitrarily high root-of-unity inverse ramification.

At infinity, with \(N=\binom n2\),
\[
 [n]_q!=q^N\left(1+\frac{n-1}{q}+\BigO(q^{-2})\right),
\]
so, on a chosen root branch,
\begin{equation}
 q=y^{1/N}-\frac{n-1}{N}+\BigO(y^{-1/N})
 =y^{1/N}-\frac2n+\BigO(y^{-1/N}).
 \label{eq:qfactorial-infinity-inverse}
\end{equation}

\chapter{q-gamma and q-digamma inverse jets}

\section{Definitions and parameter derivatives}
For \(0<q<1\),
\begin{equation}
 \Gamma_q(x)=(1-q)^{1-x}
 \frac{(q;q)_\infty}{(q^x;q)_\infty}.
 \label{eq:qgamma-def}
\end{equation}
It obeys
\(
 \Gamma_q(x+1)=[x]_q\Gamma_q(x)
\)
and \(\Gamma_q(1)=\Gamma_q(2)=1\).  The logarithmic derivative is
\begin{equation}
 \psi_q(x)=-\log(1-q)+(\log q)
 \sum_{j\ge0}\frac{q^{x+j}}{1-q^{x+j}},
 \label{eq:qdigamma}
\end{equation}
and
\begin{equation}
 \psi_q'(x)=(\log q)^2\sum_{j\ge0}
 \frac{q^{x+j}}{(1-q^{x+j})^2}>0.
 \label{eq:qtrigamma-positive}
\end{equation}
Thus \(\Gamma_q\) has the same two-branch real argument geometry as the ordinary
gamma function: its logarithm is strictly convex and has one minimum.

\section{All-order algebraic expansion at \(q=1\)}
Set \(q=\EulerE^{-\varepsilon}\).  The following coefficient functions are
polynomials in \(x\):
\begin{align}
 c_1(x)&=-\frac{(x-1)(x-2)}4, \label{eq:qgamma-c1}\\
 c_{2r}(x)&=
 \frac{B_{2r}}{2r(2r)!}
 \left(
 \frac{B_{2r+1}(x)-B_{2r+1}(1)}{2r+1}-(x-1)
 \right),\quad r\ge1. \label{eq:qgamma-c2r}
\end{align}

\begin{theorem}[\statusproved\ q-gamma expansion and odd-power cancellation]
Uniformly for \(x\) in compact subsets of the positive half-plane,
\begin{equation}
 \log\Gamma_{\EulerE^{-\varepsilon}}(x)
 \sim \log\Gamma(x)+c_1(x)\varepsilon
 +\sum_{r\ge1}c_{2r}(x)\varepsilon^{2r}.
 \label{eq:qgamma-expansion}
\end{equation}
No odd algebraic power beyond \(\varepsilon^1\) occurs.  In particular,
\begin{align}
 c_2(x)&=\frac{(x-2)(x-1)(2x+3)}{144}, \label{eq:qgamma-c2}\\
 c_4(x)&=-\frac{(x-2)(x-1)(6x^3+3x^2+7x+15)}{86400}.
 \label{eq:qgamma-c4}
\end{align}
\end{theorem}
\begin{proof}
Subtract \(\log\Gamma(x)\) from the logarithm of the functional equation.  By
\cref{eq:qnumber-q1},
\begin{equation}
 \log\frac{[x]_{\EulerE^{-\varepsilon}}}{x}
 =-\frac{x-1}{2}\varepsilon
 +\sum_{r\ge1}\frac{B_{2r}(x^{2r}-1)}{2r(2r)!}
 \varepsilon^{2r}.
 \label{eq:qgamma-difference-source}
\end{equation}
If \(C_m(x)\) is the coefficient of \(\varepsilon^m\) in the left side of
\cref{eq:qgamma-expansion}, then
\[
 C_m(x+1)-C_m(x)=[\varepsilon^m]
 \log\frac{[x]_{\EulerE^{-\varepsilon}}}{x},
 \qquad C_m(1)=0.
\]
The Bernoulli identity
\(B_{r+1}(x+1)-B_{r+1}(x)=(r+1)x^r\) gives
\cref{eq:qgamma-c1,eq:qgamma-c2r}.  The product asymptotics of
Moak and McIntosh \cite{Moak1984,McIntosh1999} select these polynomial
solutions and justify the asymptotic expansion; the difference source has no
odd powers beyond one, so neither does the selected solution.
\end{proof}

\section{Inverting the base}
Fix \(x\notin\{1,2\}\), and put
\[
 v=\log\frac{y}{\Gamma(x)}.
\]
Since \(c_1(x)\ne0\), the branch \(q\to1\) is ordinary:
\begin{equation}
 \varepsilon=\frac{v}{c_1}
 -\frac{c_2}{c_1^3}v^2
 +\frac{2c_2^2}{c_1^5}v^3+\BigO(v^4),
 \qquad q=\EulerE^{-\varepsilon}.
 \label{eq:qgamma-base-inverse}
\end{equation}
At \(x=1\) and \(x=2\), \(\Gamma_q(x)\equiv1\), so no local base inverse
exists.  These are complete, not merely first-order, degeneracies.

\section{Inverting the argument}
Let \(x_0\) be a chosen branch of the ordinary inverse gamma function,
\(\Gamma(x_0)=y\), and assume \(\psi(x_0)\ne0\).  Seek
\(x=x_0+A\varepsilon+B\varepsilon^2+\cdots\).  From
\cref{eq:qgamma-expansion},
\begin{align}
 A&=-\frac{c_1(x_0)}{\psi(x_0)}, \label{eq:qgamma-x-A}\\
 B&=-\frac{
 \frac12\psi'(x_0)A^2+c_1'(x_0)A+c_2(x_0)}{\psi(x_0)}.
 \label{eq:qgamma-x-B}
\end{align}
At the unique positive minimum, \(\psi(x_0)=0\) and
\(\psi'(x_0)>0\); the two inverse-gamma branches merge with square-root
scaling.  A simultaneous perturbation of the target and \(q\) is handled by
the fold-unfolding model of Part I.

\section{q-digamma argument and base inverses}
Equation \cref{eq:qtrigamma-positive} shows that
\(x\mapsto\psi_q(x)\) is strictly increasing.  Its range is
\(( -\infty,-\log(1-q))\), so every target in that interval has exactly one
positive real argument inverse.

Differentiating \cref{eq:qgamma-expansion} gives
\begin{equation}
 \psi_{\EulerE^{-\varepsilon}}(x)
 \sim\psi(x)+c_1'(x)\varepsilon+
 \sum_{r\ge1}c_{2r}'(x)\varepsilon^{2r}.
 \label{eq:qdigamma-q1}
\end{equation}
For a fixed target and an ordinary inverse-digamma point \(x_0\),
\[
 x=x_0-\frac{c_1'(x_0)}{\psi'(x_0)}\varepsilon+\BigO(\varepsilon^2).
\]
For base inversion at fixed \(x\), the linear coefficient
\[
 c_1'(x)=\frac{3-2x}{4}
\]
vanishes only at \(x=3/2\).  Since
\(c_2'(3/2)=-1/288\), this point has a square-root base inverse:
\begin{equation}
 \psi_{\EulerE^{-\varepsilon}}(3/2)-\psi(3/2)
 =-\frac{\varepsilon^2}{288}+\BigO(\varepsilon^4).
 \label{eq:qdigamma-critical}
\end{equation}
All other fixed positive \(x\) have an ordinary base inverse at \(q=1\).

\chapter{q-beta: coordinate inverses and a critical base curve}

\section{Coordinate monotonicity}
Define
\begin{equation}
 B_q(x,y)=\frac{\Gamma_q(x)\Gamma_q(y)}{\Gamma_q(x+y)},
 \qquad x,y>0.
 \label{eq:qbeta-def}
\end{equation}
Then
\begin{align}
 \partial_x\log B_q(x,y)&=\psi_q(x)-\psi_q(x+y)<0, \\
 \partial_y\log B_q(x,y)&=\psi_q(y)-\psi_q(x+y)<0.
 \label{eq:qbeta-coordinate}
\end{align}
Hence each coordinate inverse is unique on its positive real domain whenever
the target lies in the corresponding range.  Its local coefficients follow by
differentiating \cref{eq:qbeta-coordinate}; the first derivative is the
reciprocal of the displayed difference.

\section{Base expansion and the curve \(xy=1\)}
Combining three copies of \cref{eq:qgamma-expansion} gives
\begin{equation}
 \log\frac{B_{\EulerE^{-\varepsilon}}(x,y)}{B(x,y)}
 =b_1(x,y)\varepsilon+b_2(x,y)\varepsilon^2+\BigO(\varepsilon^4),
 \label{eq:qbeta-q1}
\end{equation}
where
\begin{align}
 b_1(x,y)&=\frac{xy-1}{2}, \label{eq:qbeta-b1}\\
 b_2(x,y)&=-\frac{x^2y+xy^2-xy-1}{24}.
 \label{eq:qbeta-b2}
\end{align}
There is no cubic term.

\begin{theorem}[\statusproved\ Type change of the q-beta base inverse]
At \(q=1\), the inverse in \(q\) is ordinary if \(xy\ne1\).  On the curve
\(xy=1\),
\begin{equation}
 b_2(x,1/x)=-\frac{(x-1)^2}{24x}.
 \label{eq:qbeta-critical-b2}
\end{equation}
Thus every positive point of that curve except \((1,1)\) is a quadratic fold
and has a square-root base inverse.  At \((1,1)\),
\(B_q(1,1)\equiv1\), so the base is completely unidentifiable.
\end{theorem}
\begin{proof}
Off \(xy=1\), \cref{eq:qbeta-b1} and the inverse-function theorem apply.  On
the curve substitute \(y=1/x\) into \cref{eq:qbeta-b2}.  The result is nonzero
unless \(x=1\).  The identity at \((1,1)\) follows from
\(\Gamma_q(1)=\Gamma_q(2)=1\).
\end{proof}

If \(v=\log(B_q/B)\), then off the critical curve
\[
 \varepsilon=\frac{v}{b_1}-\frac{b_2}{b_1^3}v^2+\BigO(v^3).
\]
On \(xy=1\), \(x\ne1\),
\[
 \varepsilon=\pm\sqrt{\frac{v}{b_2}}+\BigO(v).
\]
For the real branch \(0<q<1\), retain the sign with \(\varepsilon>0\).

\chapter{Euler q-exponentials, basic hypergeometric series, and theta products}

\section{Euler q-exponentials reduce exactly to Pochhammer inversion}
For \(0<\AbsoluteValue q<1\), use the conventions
\[
 e_q(z)=\frac1{(z;q)_\infty},
 \qquad
 E_q(z)=(-z;q)_\infty.
\]
If \(\Aq\) denotes a selected inverse branch of the infinite product in its
argument, then
\begin{equation}
 e_q^{-1}(y)=\Aq(y^{-1}),
 \qquad
 E_q^{-1}(y)=-\Aq(y).
 \label{eq:qexp-exact-inverse}
\end{equation}
These identities transport every zero, critical point, endpoint expansion,
and branch label from Part III without recomputation.  Base inversion of the
q-exponentials is likewise the base inversion of the corresponding product
with a transformed target.

Near \(z=0\), \cref{eq:infinite-a-near-zero} immediately gives
\[
 e_q^{-1}(1+s)=(1-q)s+\BigO(s^2),
 \qquad
 E_q^{-1}(1+s)=(1-q)s+\BigO(s^2),
\]
with different higher coefficients because one target is reciprocated and the
other argument is negated.

\section{A coefficientwise inverse for basic hypergeometric series}
Let a chosen basic hypergeometric function be written locally as
\[
 F(z;\bm a,\bm b,q)=1+c_1z+c_2z^2+c_3z^3+\cdots,
 \qquad c_1\ne0.
\]
Then the argument inverse is universally
\begin{equation}
 z=\frac{y-1}{c_1}-\frac{c_2}{c_1^3}(y-1)^2
 +\left(\frac{2c_2^2}{c_1^5}-\frac{c_3}{c_1^4}\right)(y-1)^3+\cdots.
 \label{eq:basic-hypergeometric-zinverse}
\end{equation}
Every \(c_n\) is a quotient of finite q-Pochhammer products, so its derivatives
in any numerator or denominator parameter reduce to finite sums of the form
\cref{eq:a-log-derivatives}.  This gives mixed inverse jets in the argument,
base, and all shape parameters.

If a numerator parameter crosses \(a_i=q^{-N}\), the series truncates.  The
argument inverse then becomes an algebraic covering of degree at most
\(N\).  If simultaneously a denominator parameter approaches a cancelling
value, one must first cancel common q-Pochhammer factors; otherwise a removable
singularity is incorrectly reported as an inverse branch point.

\section{Parameter inversion near a truncation locus}
Write \(a_i=q^{-N}(1+\eta)\).  The factor \((a_i;q)_n\) contains one small term
when \(n>N\), so the tail of the hypergeometric series is \(\BigO(\eta)\).
At \(\eta=0\), the function is a polynomial \(F_N(z)\).  If
\(\partial_\eta F\ne0\), \(\eta\) has an ordinary inverse at fixed \(z\).  If
that derivative vanishes at a root or critical point of \(F_N\), the joint
Newton polygon in \((z-z_0,\eta,y-y_0)\) determines the fractional scaling.
This is the basic-hypergeometric analogue of the finite-product fold
unfolding.

\section{Multiplicative theta products}
Use
\begin{equation}
 \vartheta(z;q)=(z;q)_\infty(q/z;q)_\infty(q;q)_\infty.
 \label{eq:multiplicative-theta}
\end{equation}
For fixed \(0<\AbsoluteValue q<1\), its zeros \(z=q^m\), \(m\in\IntegerNumbers\), are simple.  Hence
the argument inverse at target zero is ordinary on every sheet:
\[
 z=q^m+\frac{y}{\partial_z\vartheta(q^m;q)}+\BigO(y^2).
\]
The quasi-periodicity
\(
 \vartheta(qz;q)=-z^{-1}\vartheta(z;q)
\)
transports the derivative normalization from one zero to all others.

Base inversion is more delicate.  Near \(q=0\), the Laurent/Fourier series gives
an ordinary or Puiseux jet according to the first nonzero q-power.  Near
\(q=1\), modular transformation produces exponential scales
\(\EulerE^{-\pi^2/\varepsilon}\), so a complete inverse is a transseries rather than
a Bernoulli power series.  The dominant exponential is inverted by a logarithm
or Lambert \(W\); the modular images determine the Stokes corrections.

\section{Other positive palindromic q-enumerators}
Whenever a finite q-analog is a nonconstant polynomial with nonnegative
palindromic coefficients, the centered theorem of Part IV applies without
change.  Examples include Poincare polynomials of finite Coxeter groups,
q-multinomials, many Gaussian quotients, and selected q-Catalan and plane
partition enumerators under their palindromic normalizations.  Objects lacking
palindromicity still have ordinary local inverse jets, but lose the exact
reciprocal two-branch geometry.

\part{Fabius--Rvachev synthesis, algorithms, and frontiers}

\chapter{An effective geometric base for the Rvachev infinite sinc product}

\section{A one-parameter compactly supported convolution family}
Let \(\SincRad z=\sin z/z\), and for \(0<q<1\) define
\begin{equation}
 \widehat U_q(t)=\prod_{j\ge0}
 \SincRad\bigl((1-q)q^jt\bigr).
 \label{eq:geometric-sinc-family}
\end{equation}
Probabilistically, this is the characteristic function of the sum of
independent centered uniform random variables whose half-widths are
\((1-q)q^j\).  Their widths sum to one, so the corresponding density is
supported on \([-1,1]\).  At \(q=1/2\),
\[
 \widehat U_{1/2}(t)=\prod_{r\ge1}\SincRad(t/2^r),
\]
which is the standard infinite sinc product underlying the Rvachev up-function,
up to Fourier-transform normalization.  Its cumulative form is correspondingly
linked to the Fabius function in the conventions developed by the source
corpus.

\section{Cumulants and a globally identifiable base}
Euler's sine product gives
\[
 \log\SincRad z
 =-\sum_{m\ge1}\frac{\zeta(2m)}{m\pi^{2m}}z^{2m}.
\]
Therefore
\begin{equation}
 \log\widehat U_q(t)=
 -\sum_{m\ge1}\frac{\zeta(2m)}{m\pi^{2m}}
 R_m(q)t^{2m},
 \qquad
 R_m(q)=\frac{(1-q)^{2m}}{1-q^{2m}}.
 \label{eq:geometric-sinc-cumulants}
\end{equation}

\begin{theorem}[\statusproved\ Every even cumulant identifies the geometric
base]
For every \(m\ge1\), \(R_m\) is strictly decreasing from \(1\) to \(0\) on
\((0,1)\).  Hence any one nonzero even cumulant coefficient in
\cref{eq:geometric-sinc-cumulants} determines \(q\) uniquely.
\end{theorem}
\begin{proof}
Differentiate the logarithm:
\[
 \frac{R_m'(q)}{R_m(q)}
 =-\frac{2m}{1-q}+\frac{2mq^{2m-1}}{1-q^{2m}}.
\]
The second fraction is smaller than the first because
\(q^{2m-1}(1-q)<1-q^{2m}\), equivalently \(q^{2m-1}<1\).  The endpoint limits
are immediate.
\end{proof}

At the dyadic point,
\begin{equation}
 R_m(1/2)=\frac1{2^{2m}-1},
 \qquad
 R_m'(1/2)=-\frac{4m(2^{2m}-2)}{(2^{2m}-1)^2}.
 \label{eq:dyadic-Rm}
\end{equation}
Thus if an observed normalized cumulant ratio is
\(R_m=1/(2^{2m}-1)+\delta\), its effective-base inverse is
\begin{equation}
 q=\frac12-
 \frac{(2^{2m}-1)^2}{4m(2^{2m}-2)}\delta+\BigO(\delta^2).
 \label{eq:effective-dyadic-base}
\end{equation}
This gives a direct inverse bridge between Fabius--Rvachev moment data and the
q-analog machinery of the preceding parts.

\section{Overdetermined consistency and deformations}
A genuine geometric convolution must return the same \(q\) from every even
cumulant.  Disagreement between the inverses of \(R_1,R_2,\ldots\) therefore
measures departure from exact geometric scaling.  With an additional global
scale \(s\), replace each width by \(s(1-q)q^j\); the \(2m\)-th logarithmic
coefficient becomes \(s^{2m}R_m(q)\).  Two independent cumulants generically
recover \((s,q)\), while further cumulants test the model.

The same idea applies to finite geometric combs.  Their Fourier products are
finite q-products in the squared frequency, and their interpolation weights
contain Gaussian coefficients.  Inverting an observed comb response for its
ratio \(q\) is therefore a simultaneous inverse problem in the sense of
\cref{eq:good}, with the dyadic comb \(q=1/2\) as an ordinary interior base
point rather than a singular endpoint.

\chapter{Numerical realization of the inverse atlas}

\section{Coordinate selection before iteration}
A reliable solver first changes coordinates so that the desired branch is
regular:
\begin{center}
\begin{tabularx}{0.94\textwidth}{@{}lXX@{}}
\toprule
Regime & Forward coordinate & Inverse coordinate \\
\midrule
\(q\approx0\) & \(q\) or \(q^{1/s}\) & ordinary/Puiseux series \\
\(q\approx1\) & \(\varepsilon=-\PrincipalLogarithm q\) & logarithmic target \(w=\PrincipalLogarithm(y/y_0)\) \\
\(q\approx-1\) & \(h=q+1\) & \(h=(y-y_0)^{1/m}v\) \\
root \(\zeta\) & \(h=q-\zeta\) & cyclotomic multiplicity coordinate \\
\(q\approx\infty\) & \(Q=q^{-1}\) & \(Y=y^{1/D}\), then Laurent reversion \\
fixed-target infinite product & \((a,\varepsilon)\) & weighted degree \(a\sim\varepsilon\) \\
\bottomrule
\end{tabularx}
\end{center}
Using Newton iteration in the original variable near a fold loses half the
digits and can jump sheets.  Iteration in the desingularized coordinate is both
better conditioned and branch explicit.

\section{Predictor--corrector continuation}
For a regular branch \(F(p,\lambda)=y\), differentiate along a continuation
path \((y(s),\lambda(s))\):
\begin{equation}
 \frac{dp}{ds}=
 \frac{y'(s)-F_\lambda\lambda'(s)}{F_p}.
 \label{eq:continuation-ode}
\end{equation}
An Euler or higher-order predictor from \cref{eq:continuation-ode} is followed
by a Newton corrector.  Near \(F_p=0\), switch to pseudo-arclength continuation
of the curve \((p,y)\); this traverses a fold without choosing an artificial
vertical graph.  A branch label consists of the starting germ plus the ordered
list of discriminant cuts crossed.

\section{Exact root-of-unity preprocessing}
For polynomial q-analogs, factor cyclotomically before numerical work.
\begin{enumerate}
\item Compute the exponent of every \(\Phi_d\) by floor formulas rather than by
floating-point factorization.
\item Divide out the exact cyclotomic factor and evaluate the nonzero cofactor
at \(\zeta\).
\item Use the multiplicity as the Puiseux denominator and the cofactor to
normalize the leading root.
\end{enumerate}
This prevents a cluster of nearly coincident numerical roots from being
mistaken for distinct simple branches.

\section{Certified local inverses}
Suppose a truncated inverse jet \(p_M(y)\) is available on a disk.  Compute the
residual
\[
 R_M(y)=F(p_M(y))-y.
\]
If an interval or ball enclosure proves \(F_p\) stays away from zero, then a
Newton--Kantorovich bound turns \(R_M\) into an error bound for the true branch.
At a ramification point, apply the same argument to the desingularized equation
in \(v=(p-p_0)/(y-y_0)^{1/m}\).  Rouch\'e's theorem can additionally certify
the number of branches inside a contour.

\section{Reproducible experiments in the archive}
The accompanying program
\texttt{inverse\_q\_analog\_experiments.py} uses exact SymPy algebra and
high-precision mpmath arithmetic.  It verifies the following independent
claims.
\begin{enumerate}
\item The parameter \cref{eq:a-star-minus-one}, the logarithmic second
derivatives \cref{eq:L2-even,eq:L2-odd}, and the exact cubic
\cref{eq:exact-cubic-minus-one}.
\item Cyclotomic carry counts, Gaussian square-freeness in sample families,
the q-factorial coefficient \(m!\) at \(q=-1\), and
\cref{eq:gaussian-minus-one-derivative}.
\item The five-term double-scaling inverse \cref{eq:double-scaling-five} at two
targets and three decreasing values of \(\varepsilon\).
\item The q-gamma coefficients \cref{eq:qgamma-c1,eq:qgamma-c2,eq:qgamma-c4}
and the absence of a cubic algebraic term.
\item The centered Gaussian Puiseux expansion for \((n,k)=(8,3)\), where
\[
 \kappa_2=15,
 \qquad \kappa_4=-111,
 \qquad \kappa_6=2520.
\]
\item The q-beta coefficients and the reduction
\(b_2(x,1/x)=-(x-1)^2/(24x)\).
\end{enumerate}

A representative double-scaling check is
\begin{center}
\begin{tabular}{@{}cccc@{}}
\toprule
\(y\) & \(\varepsilon\) & numerical \(a\) & five-term error \\
\midrule
\(0.3\) & \(0.05\) & \(0.0560243496070\) & \(-8.05\times10^{-9}\) \\
\(0.3\) & \(0.025\) & \(0.0288380205105\) & \(-1.29\times10^{-10}\) \\
\(0.7\) & \(0.05\) & \(0.0173299163772\) & \(-4.24\times10^{-11}\) \\
\(0.7\) & \(0.025\) & \(0.00878272387187\) & \(-6.70\times10^{-13}\) \\
\bottomrule
\end{tabular}
\end{center}
The archive includes the full text output so that no rounded table is the sole
record of a computation.

\chapter{Conjectures and research directions}

\section{Algebraic monodromy of finite base inverses}
For fixed generic \((a,y)\), the equation
\(P_n(a,q)=y\) has degree \(N=\binom n2\) in \(q\).  Its coefficients are highly
structured but possess no obvious decomposition for generic parameters.

\begin{conjecture}[Generic symmetric monodromy]
For \(n\ge3\), after removing the obvious degenerate loci
\(a\in\{0,1\}\), \(y\) critical, and lower-degree specializations, the
geometric monodromy group of \(P_n(a,q)-y\) over the \((a,y)\)-parameter space
is the full symmetric group \(S_N\).
\end{conjecture}

A computational route is to specialize to several rational pairs \((a,y)\),
compute exact Galois groups modulo good primes, and then prove that a
transposition and a long cycle occur in the generic monodromy.  The explicit
folds at \(q=-1\) naturally supply transposition candidates.

\section{Secondary discriminants and collision normal forms}
The primary discriminant eliminates \(q\) from \(P_n-y=0\) and
\(\partial_qP_n=0\).  A secondary discriminant detects equality of two distinct
critical values.  The first inverse report found a positive-base collision for
\((a;q)_5\).  The next objective is a stratified classification rather than a
list of examples.

\begin{conjecture}[Generic normal crossings of critical-value collisions]
Away from roots of unity, identically constant parameters, and higher critical
points, the secondary-discriminant hypersurface for finite q-Pochhammer base
inversion is reduced, and a generic point represents the transverse equality
of exactly two simple critical values.
\end{conjecture}

A proof would combine elimination theory with transversality of the critical
value map.  Exceptional strata should correspond to the vanishing of explicit
higher resultants.

\section{Resurgence of \(q\to1\) inverse jets}
The Bernoulli coefficients in \cref{eq:infinite-q1-forward} grow factorially,
and modular/theta analogs contain exponentially small scales invisible to the
power series.

\begin{conjecture}[Sectorial Borel summability]
For fixed \(a\) away from the polylogarithmic branch locus, the formal base
inverse \cref{eq:infinite-base-q1-inverse} and the corresponding q-gamma and
q-beta inverse jets are Gevrey-1 and Borel summable in every non-Stokes sector.
Their Stokes discontinuities are generated by the nearest dilogarithmic or
modular saddle actions.
\end{conjecture}

The technically promising route is to begin from integral representations of
q-shifted factorials, prove uniform alien-derivative estimates for the forward
transseries, and use closure of resurgent functions under local composition and
reversion.

\section{Global q-beta critical surfaces}
The local unit-base transition is the simple curve \(xy=1\), but numerical
sampling shows that the sign of \(\partial_qB_q(x,y)\) can change inside
\((0,1)\) for several positive parameter pairs.  Thus the global critical set
is richer than a continuation of \(xy=1\).

\begin{researchproblem}
Determine the real analytic surface
\[
 \mathcal C=\{(x,y,q)>0:0<q<1,
 \;\partial_qB_q(x,y)=0\},
\]
its number of sheets over the \((x,y)\)-plane, and the cusp or Maxwell loci of
its critical values.  Establish parameter regions with zero, one, or multiple
positive-base inverse branches.
\end{researchproblem}

Lambert-series differentiation of \cref{eq:qbeta-def} gives a rapidly
convergent exact equation for \(\mathcal C\); interval continuation can turn the
initial numerical atlas into proofs.

\section{Root-of-unity universality for multinomial inverses}
The carry theorem fixes the ramification order but not the full unfolding.
Allowing the integer parts to move through q-gamma interpolation produces
continuous parameters near the cyclotomic zero.

\begin{conjecture}[Generic \(A_{m-1}\) unfolding]
At a primitive \(d\)-th root with q-multinomial carry multiplicity \(m\), a
generic \((m-1)\)-parameter continuous deformation of the parts gives a
versal unfolding of the singularity \(y=c(q-\zeta)^m\), equivalent to the
\(A_{m-1}\) normal form.
\end{conjecture}

This would connect elementary floor carries to the singularity theory of
inverse coverings and would predict all local branch-merger diagrams.

\section{Coefficient arithmetic of the double-scaling inverse}
The polynomials \(A_m(L)\) in \cref{eq:double-scaling-five} are generated by
Bernoulli numbers and partitioned compositions.  Their denominators and signs
show nontrivial structure.

\begin{researchproblem}
Find a closed formula for \(A_m(L)\) in partial Bell polynomials, determine the
minimal denominator as a function of \(m\), and study the zero distribution of
\(A_m(L)/L\).  Decide whether a natural Bernoulli rescaling produces integer
coefficients or a combinatorial interpretation.
\end{researchproblem}

\section{Fabius--Rvachev inverse consistency}
The family \cref{eq:geometric-sinc-family} provides a direct base estimator
from each even cumulant.  For exact geometric scaling all estimators coincide.
A broader inverse problem allows the widths to be
\(w_j=(1-q)q^j(1+\epsilon_j)\) with a small structured perturbation.

\begin{researchproblem}
Recover finitely many deformation modes \((q,\epsilon_0,\ldots,\epsilon_s)\)
from a finite moment vector, determine the exact Jacobian rank at the dyadic
point, and derive minimax conditioning bounds.  Relate loss of rank to the
flatness and endpoint asymptotics of the associated Fabius/Rvachev functions.
\end{researchproblem}

Good's multivariate formula supplies formal inverse coefficients; the remaining
work is global identifiability and analytic error control.

\section{Formal verification roadmap}
The algebraic core is well suited to Lean formalization in the following order.
\begin{enumerate}
\item finite formal-series reversion and the triangular recurrence;
\item cyclotomic exponents for Gaussian coefficients, q-multinomials, and
q-factorials;
\item the carry theorem and square-free Gaussian corollary;
\item exact \(q=-1\) identities, including \cref{eq:exact-cubic-minus-one} and
\cref{eq:gaussian-minus-one-derivative};
\item positivity and the centered-palindromic monotonicity theorem;
\item finite-order coefficient identities for the Bernoulli jets;
\item analytic convergence, asymptotic remainder estimates, and sectorial
continuation.
\end{enumerate}
The first five items require only polynomial algebra, finite sums, and real
monotonicity.  They can therefore be separated cleanly from the much heavier
complex asymptotic infrastructure.

\part{Reference appendices}

\chapter{Inverse-regime tables}

\section{Finite q-Pochhammer product}
\begin{longtable}{@{}p{0.13\textwidth}p{0.18\textwidth}p{0.16\textwidth}p{0.43\textwidth}@{}}
\caption{Local inverse types for \((a;q)_n\).}\label{tab:p-atlas}\\
\toprule
Recovered variable & Base point & Type & Leading coordinate/criterion \\
\midrule
\endfirsthead
\toprule
Recovered variable & Base point & Type & Leading coordinate/criterion \\
\midrule
\endhead
\(a\) & \(a=0,y=1\) & Taylor & derivative \(-[n]_q\); singular when
\([n]_q=0\) \\
\(a\) & \(a=q^{-j},y=0\) & Taylor & exact derivative
\cref{eq:derivative-geometric-zero} \\
\(a\) & \(a=\infty\) & \(n\)-Puiseux &
\(a\sim(y/((-1)^nq^{\binom n2}))^{1/n}\) \\
\(q\) & \(q=0\) & Taylor & normalized target \(r\),
\cref{eq:subset-small-q} \\
\(q\) & \(q=1\) & Taylor & \(\varepsilon=-\PrincipalLogarithm q\), logarithmic target
\(w\) \\
\(q\) & \(q=-1,a=a_*\) & square-root & fold for \(n\ge4\) \\
\(q\) & \(q=-1,n=3,a=-1/3\) & cubic root & exact
\cref{eq:exact-cubic-minus-one} \\
\(q\) & root \(\zeta\), target zero & \(m\)-Puiseux & multiplicity count
\cref{eq:p-root-unity-multiplicity} \\
\(q\) & \(q=\infty\) & \(\binom n2\)-Puiseux & reciprocal factorization
\cref{eq:p-reciprocal-factorization} \\
\(\alpha\) & regular complex order & Taylor & derivative
\cref{eq:alpha-derivatives} \\
\bottomrule
\end{longtable}

\section{Gaussian and related polynomial q-analogs}
\begin{longtable}{@{}p{0.18\textwidth}p{0.17\textwidth}p{0.17\textwidth}p{0.39\textwidth}@{}}
\caption{Local inverse types for Gaussian and related polynomials.}\label{tab:g-atlas}\\
\toprule
Object & Base point & Type & Leading mechanism \\
\midrule
\endfirsthead
\toprule
Object & Base point & Type & Leading mechanism \\
\midrule
\endhead
Gaussian & \(q=0,y=1\) & Taylor & restricted partitions \\
Gaussian & \(q=1\), uncentered & Taylor & deterministic drift \(-d/2\) \\
Gaussian & \(q=1\), centered & square-root & variance/cumulants \\
Gaussian & primitive root, target zero & Taylor & square-free cyclotomic factor \\
Gaussian & \(q=-1\), zero case & Taylor & exact derivative
\cref{eq:gaussian-minus-one-derivative} \\
Gaussian & \(q=\infty\) & \(d\)-Puiseux & reciprocity \\
q-multinomial & primitive root & \(e_d\)-Puiseux & residue-carry count \\
q-factorial & primitive root & \(\lfloor n/d\rfloor\)-Puiseux & multiples of
\(d\) \\
positive palindromic polynomial & centered \(q=1\) & square-root & symmetric
moment-generating function \\
\bottomrule
\end{longtable}

\chapter{Coefficient formula compendium}

\section{Scalar ordinary reversion}
For \(w=c_1z+c_2z^2+c_3z^3+c_4z^4+\cdots\),
\begin{align*}
 z={}&\frac{w}{c_1}-\frac{c_2}{c_1^3}w^2
 +\frac{2c_2^2-c_1c_3}{c_1^5}w^3 \\
 &+\frac{-5c_2^3+5c_1c_2c_3-c_1^2c_4}{c_1^7}w^4
 +\BigO(w^5).
\end{align*}
For implementation, the recurrence in Part I avoids expression swell from
closed Bell-polynomial formulas.

\section{Fold inversion}
If
\[
 y-y_0=A z^2+Bz^3+Cz^4+\cdots,
 \qquad A\ne0,
\]
and \(\xi=((y-y_0)/A)^{1/2}\), then
\[
 z=\pm\xi-\frac{B}{2A}\xi^2
 \pm\left(\frac{5B^2}{8A^2}-\frac{C}{2A}\right)\xi^3
 +\BigO(\xi^4).
\]
The signs on odd powers follow the selected square-root sheet.

\section{Centered cumulant inversion}
For \(u=\sum_{r\ge1}\kappa_{2r}\varepsilon^{2r}/(2r)!\), first revert the
ordinary series in \(z=\varepsilon^2\), then take a square root.  This two-stage
method is algebraically more stable than direct Puiseux coefficient matching.

\section{Reciprocal polynomial inversion}
If
\[
 F(q)=Cq^D(1+c_1q^{-1}+c_2q^{-2}+c_3q^{-3}+\cdots),
 \qquad Y=(y/C)^{1/D},
\]
then
\begin{align*}
 q={}&Y-\frac{c_1}{D}
 +\left(\frac{c_1^2(D-1)}{2D^2}-\frac{c_2}{D}\right)Y^{-1} \\
 &+\left(
 -\frac{c_1^3(D-1)(2D-1)}{3D^3}
 +\frac{2c_1c_2(D-1)}{D^2}
 -\frac{c_3}{D}
 \right)Y^{-2}+\BigO(Y^{-3}).
\end{align*}
This single template yields
\cref{eq:p-large-q-inverse,eq:gaussian-large-q,eq:qfactorial-infinity-inverse}.

\chapter{Reproduction notes and source map}

\section{Archive contents}
The companion archive contains:
\begin{description}
\item[\texttt{inverse\_q\_analog\_jet\_atlas.tex}] the complete LaTeX source;
\item[\texttt{inverse\_q\_analog\_jet\_atlas.pdf}] the compiled report;
\item[\texttt{inverse\_q\_analog\_experiments.py}] exact and high-precision
checks with detailed comments;
\item[\texttt{numerical\_checks.txt}] the deterministic output generated at
80 decimal digits;
\item[\texttt{README.txt}] build and reproduction commands.
\end{description}

\section{Principal audited TeX sources}
The live source tree at the pinned commit contains, among its TeX material, the
following principal documents:
\begin{itemize}
\item \texttt{Exponents\_and\_q\_Series\_Frontiers/
Exponents\_and\_q\_Series\_Frontiers.tex};
\item \texttt{Fabius\_Flat\_Parameter\_Response\_Dynamics/
fabius\_frontier\_report.tex};
\item \texttt{Fabius\_Rvachev\_Frontier\_Report/
fabius\_frontier\_report.tex} and its generated
\texttt{numerical\_results.tex};
\item \texttt{inverse\_q\_analogs\_report/
inverse\_q\_analogs\_report.tex};
\item \texttt{q\_pochhammer\_q\_binomial\_expansions\_report/
q\_analog\_expansions\_report.tex};
\item \texttt{q\_pochhammer\_q\_binomial\_monograph/
q\_pochhammer\_q\_binomial\_monograph.tex};
\item small generated TeX table fragments under the synthesis assets.
\end{itemize}
The forward identities were cross-checked against the proof-oriented monograph;
the new inverse statements were then derived independently and verified where
practical by the archive script.

\section{Build commands}
From the archive directory, a standard TeX installation can reproduce the PDF
with
\begin{lstlisting}[basicstyle=\ttfamily\small,frame=single]
latexmk -pdf -interaction=nonstopmode -halt-on-error \
  inverse_q_analog_jet_atlas.tex
\end{lstlisting}
The numerical audit is regenerated with
\begin{lstlisting}[style=pythonstyle]
python inverse_q_analog_experiments.py --digits 80 \
  --output numerical_checks.txt
\end{lstlisting}
The script requires only SymPy and mpmath.

\backmatter

\begin{thebibliography}{99}

\bibitem{Andrews1976}
G. E. Andrews,
\emph{The Theory of Partitions},
Encyclopedia of Mathematics and its Applications, Vol. 2,
Addison--Wesley, 1976; reprinted by Cambridge University Press.

\bibitem{AndrewsAskeyRoy}
G. E. Andrews, R. Askey, and R. Roy,
\emph{Special Functions},
Encyclopedia of Mathematics and its Applications, Vol. 71,
Cambridge University Press, 1999.

\bibitem{Comtet}
L. Comtet,
\emph{Advanced Combinatorics: The Art of Finite and Infinite Expansions},
D. Reidel, 1974.

\bibitem{DLMF}
NIST Digital Library of Mathematical Functions,
Chapters 5 and 17: Gamma and q-hypergeometric functions,
\url{https://dlmf.nist.gov/}, accessed August 2026.

\bibitem{Fabius1966}
J. Fabius,
A probabilistic example of a nowhere analytic \(C^\infty\)-function,
\emph{Zeitschrift f\"ur Wahrscheinlichkeitstheorie und Verwandte Gebiete}
\textbf{5} (1966), 173--174.

\bibitem{FlajoletSedgewick}
P. Flajolet and R. Sedgewick,
\emph{Analytic Combinatorics},
Cambridge University Press, 2009.

\bibitem{GasperRahman}
G. Gasper and M. Rahman,
\emph{Basic Hypergeometric Series}, second edition,
Encyclopedia of Mathematics and its Applications, Vol. 96,
Cambridge University Press, 2004.

\bibitem{Good1960}
I. J. Good,
Generalizations to several variables of Lagrange's expansion, with applications
to stochastic processes,
\emph{Proceedings of the Cambridge Philosophical Society}
\textbf{56} (1960), 367--380.

\bibitem{McIntosh1999}
R. J. McIntosh,
Some asymptotic formulae for q-shifted factorials,
\emph{Ramanujan Journal} \textbf{3} (1999), 205--214,
DOI: 10.1023/A:1006949508631.

\bibitem{Moak1984}
D. S. Moak,
The q-analogue of Stirling's formula,
\emph{Rocky Mountain Journal of Mathematics} \textbf{14} (1984), 403--413.

\bibitem{Olver}
F. W. J. Olver,
\emph{Asymptotics and Special Functions},
Academic Press, 1974; A K Peters reprint, 1997.

\bibitem{StanleyEC1}
R. P. Stanley,
\emph{Enumerative Combinatorics}, Vol. 1, second edition,
Cambridge University Press, 2012.

\bibitem{ReshetnikovMonograph}
V. Reshetnikov,
\emph{q-Pochhammer Symbols and q-Binomial Coefficients},
ProveIt repository, file
\path{Analysis/FabiusFunction/docs/semi-formalized-research-frontiers/
drafts/exponents-and-q-series/q_pochhammer_q_binomial_monograph/
q_pochhammer_q_binomial_monograph.tex},
commit \nolinkurl{1cea73234a363ddbc392816f6babb5a57920e984}, 2026.

\bibitem{ReshetnikovExpansions}
V. Reshetnikov,
\emph{Expansions of q-Pochhammer Symbols, Gaussian Coefficients, and Related
q-Analogs},
ProveIt repository, file
\texttt{q\_pochhammer\_q\_binomial\_expansions\_report/
q\_analog\_expansions\_report.tex}, same commit, 2026.

\bibitem{ReshetnikovInverse}
V. Reshetnikov,
\emph{Inverse q-Analogs: Branch Geometry, Asymptotic Inversion, and
Computation},
ProveIt repository, file
\texttt{inverse\_q\_analogs\_report/inverse\_q\_analogs\_report.tex},
same commit, 2026.

\bibitem{ReshetnikovSynthesis}
V. Reshetnikov,
\emph{Exponents and q-Series Frontiers},
ProveIt repository, file
\texttt{Exponents\_and\_q\_Series\_Frontiers/
Exponents\_and\_q\_Series\_Frontiers.tex}, same commit, 2026.

\end{thebibliography}

\end{document}
